\documentclass[11pt,a4paper]{article}
\usepackage[T1]{fontenc}
\usepackage[utf8]{inputenc}
\usepackage{amsmath,amssymb,amsthm,mathtools}
\usepackage[margin=25mm]{geometry}
\usepackage{longtable,array,booktabs}
\usepackage[colorlinks=true,linkcolor=blue,urlcolor=blue,citecolor=blue]{hyperref}
\usepackage{microtype}
\newtheorem{theorem}{Theorem}[section]
\newtheorem{lemma}[theorem]{Lemma}
\newtheorem{proposition}[theorem]{Proposition}
\newtheorem{corollary}[theorem]{Corollary}
\theoremstyle{definition}\newtheorem{definition}[theorem]{Definition}
\theoremstyle{remark}\newtheorem{remark}[theorem]{Remark}
\newcommand{\cat}{\mathbin{\Vert}}
\newcommand{\ints}[2]{[#1,#2]_{\mathbb Z}}
\newcommand{\ES}{\operatorname{ES}}
\newcommand{\Hi}{\operatorname{Hi}}
\newcommand{\Lo}{\operatorname{Lo}}
\newcommand{\mult}{\mathrel{=_{\mathrm{mult}}}}
\newcommand{\sh}[2]{(#1+#2)}
\title{Zero-rotatability of spiders with even equal long arms and arbitrary center leaves}
\author{Beedbyte}
\date{10 October 2026}
\begin{document}
\maketitle
\begin{abstract}
For every even integer $k\ge2$, every $n\ge2$, every $m\ge0$, and every
vertex of the actual named spider $S(k^n,1^m)$, there is a graceful
labeling taking zero at that vertex. Different requests use separate
labelings. Elementary constructions cover lengths two and four, with
the entire family with arm length four explicitly derived from older
subdivision and alpha-identification operations. The larger-length proof
combines balanced midpoint cores with two infinite shell families;
literal certificates and exact transfers cover lengths six through
thirty-four. All original center leaves and the two-arm boundary are
retained. A separate internal Lean source check freshly compiles all 186
project modules and verifies the exact unconditional named-graph
statement. Historical priority is unresolved.
\end{abstract}
\noindent\textbf{Authorship and verification.} Jordi Gartner owns and
publishes the work through Beedbyte. Mathematical scope and
computational/formal verification are distinguished in
Section~\ref{sec:verification}. No claim of worldwide novelty or
historical priority is made.
\tableofcontents

\section{Statement, named graph, and context}
\label{sec:statement}
Write $G=S(k^n,1^m)$ for the following explicitly named tree. Its vertices
are a designated center $c$, long-arm vertices $v_{a,d}$ with
$0\le a<n$, $1\le d\le k$, and unit leaves $u_j$ with $0\le j<m$.
Its edges are $cv_{a,1}$, $v_{a,d}v_{a,d+1}$, and $cu_j$.
There are $N=nk+m$ edges. Depth is distance from this designated center.
When $n=2,m=0$, the graph is a path, but the designated midpoint and both
arm names remain part of the construction.

A graceful labeling is a bijection $f:V(G)\to\ints0N$ such that the
multiset of absolute edge differences is $\ints1N$. A graph is
\emph{zero-rotatable} if every vertex can receive zero in some graceful
labeling.
\begin{theorem}\label{thm:main}
For every even $k\ge2$, $n\ge2$, $m\ge0$, and every named vertex $v$ of
$S(k^n,1^m)$, there is a graceful labeling $f$ with $f(v)=0$.
\end{theorem}
The existential labeling depends on $v$. The theorem asserts neither
simultaneous zero labels nor a result for arbitrary mixed long-arm lengths.
The threshold is sufficient; no impossibility below it is asserted.

Patterson's thesis~\cite[Section 5.4, Conjecture 5.4.4]{patterson}
proposes zero-rotatability for spiders outside the family $S(3,1,1,\ldots)$.
Rofa~\cite[Introduction, pp.~3--4, and Summary]{rofa} explicitly defines
symmetric spiders by equal leg lengths and conjectures zero-rotatability
for them. Here $m=0$ lies in that symmetric subclass; $m>0$
does not. Our statement addresses a particular family in these broader
directions and does not settle either general conjecture.

Balanced permutation concatenation and insertion have older precedents.
The proof of Hicks--Ollis--Schmitt's Lemma 4.5~\cite{hos} gives a
translated graceful-prefix/bipartite-suffix concatenation. Ollis's
Lemma 5.5 and Theorem 5.6~\cite{ollis} supply balanced endpoint freedom
(with Lemma 5.5 credited there to Kotzig, reference~[21]) and an insertion construction.
Those results concern their stated permutation endpoints and joins;
they do not, by themselves, prescribe the interior zero positions needed
here. All inventories and position claims used below are proved explicitly.

% V2_CONTEXT_BEGIN known_overlaps

\paragraph{Known subfamilies and mechanisms.}
The full cases $n=2,m=0$ and $n=2,m=1$ are already covered by path
zero-rotatability and the central-pendant theorem of Luiz, Campos and
Richter~\cite[Lemma 4 and Theorem 14]{lcr}. In the first case the
designated center is path index $k$, and an arm depth $d$ is index
$k-d$ or $k+d$; in the second case the original unit leaf is attached
at that midpoint. Rofa's Corollary 2 already supplies the symmetric
$m=0$ boundary targets $c$ and depths $1,k-1,k$.
Earlier subdivision and matching machinery gives further complete
subfamilies. Uniform doubling using Krop's construction~\cite[Theorem
2.4 and generalized $\Delta_{+1}$ construction]{krop}, starting from
the adjacent extreme orientations of the length-two symmetric base,
gives the $m=0$ power-of-two lengths $k\ge32$ as a derived overlap.
For $m=1$, matching $c$ to the unit leaf and depths $(2j-1,2j)$ on each
arm gives contree $S((k/2)^n)$; the Broersma--Hoede lift, as stated in
\cite[Lemma 3 and Theorem 9]{lcr}, transfers complete zero-rotatability
from that contree. It covers the same power-of-two range. With $m\ge2$
there is no perfect matching because two unit leaves share their only
neighbor. These are deductions from the older operations, not a claim
that Krop's ordinary gracefulness statement itself prescribes all zeros.

Cattell's path-label work~\cite{cattell} and its precise restatements in
\cite[Lemma 5]{lcr} and \cite[Lemma 2]{sz} concern ordinary
single-vertex label freedom or the specified alpha endpoint/zero cases.
Separate choices of an endpoint label and an internal zero do not imply
their simultaneous occurrence in one labeling. The complete original
Cattell proof was not obtained, so stronger consequences of its
$\pi$-representations remain an open comparison rather than being excluded.

Dumitru--Nacu~\cite[Theorems 3.3--3.5 and Corollary 3.6]{dn} pack
disjoint self-matched arms at hub label one, including a zero-containing
arm, and make unused labels hub leaves. This gives ordinary gracefulness
for each fixed nontrivial arm profile with sufficiently many leaves.
Their displayed closure template places zero near the chosen tip;
the published statement does not quantify every requested interior
vertex. Further full-zero consequences can be derived under restrictive
leaf-count/congruence hypotheses, but are not an unrestricted all-$m$
theorem quoted from that paper.

Thus substantial subfamilies and generic operations are prior work.
Complete prior subsumption of the present all-$n,m$, every-vertex
statement remains unresolved. Full original proofs of Cattell,
Abrham--Kotzig/Abrham and the classical Rosa/Kotzig/Huang--Kotzig--Rosa
constructions were not all accessible; precise later restatements were
used where identified. No unsuccessful search or access gap establishes
novelty.
% V2_CONTEXT_END known_overlaps

\section{Word notation and whole-graph complementation}
Words are indexed from zero. Concatenation is $\cat$; $a+W$ means
adding $a$ to every entry of $W$. An integer interval denotes a multiset
with each member once; $\mult$ denotes multiset equality. Empty intervals
and empty concatenations contribute no entries. Put
\[
 \ES(W)=(W_0+W_1,\ldots,W_{\ell-2}+W_{\ell-1}).
\]
Every inventory in this article is a multiset inventory, so a repeated
sum cannot be hidden by set notation.

\begin{lemma}[Complement]\label{lem:complement}
If $f$ is graceful on a graph with $N$ edges, then $\bar f(v)=N-f(v)$
for \emph{every} vertex is graceful. In particular, a vertex labeled $N$
can be made zero.
\end{lemma}
\begin{proof}
The map $x\mapsto N-x$ permutes $\ints0N$, and
$|(N-f(u))-(N-f(v))|=|f(u)-f(v)|$ on every edge.
\end{proof}
In particular a zero \emph{high offset} below means a full graph label
$N$, not a literal zero. Complementation must also change the center,
all other arms, and every original unit leaf. Permuting long-arm names is
a graph automorphism and permits any selected arm.

\section{Center-zero residual and elementary targets}
\label{sec:centerzero}
Let $k=2R$, $h\ge0$, and $m\ge0$. On $S(k^h,1^m)$ define
\begin{align*}
 b(c)&=0,& b(u_j)&=hk+j+1,\\
 b(v_{i,2t})&=ik+t\quad(1\le t\le R),&
 b(v_{i,2t+1})&=(h-i)k-t\quad(0\le t<R).
\end{align*}
\begin{lemma}\label{lem:centerzero}
This labeling is graceful for every $h\ge0,m\ge0$.
\end{lemma}
\begin{proof}
The even-depth entries fill the first half of each block of $k$ positive
integers, and the odd-depth entries fill the second halves in reversed
arm order. Together they are $\ints1{hk}$. Unit leaves finish the label
interval. The root edges have weights $k,2k,\ldots,hk$.
The internal edges on arm $i$ have weights
$|(h-2i)k-s|$, $1\le s<k$. If $j=h-2i>0$, this fills the open
block between $(j-1)k$ and $jk$; if $j\le0$, it fills the open block
between $(-j)k$ and $(1-j)k$. The block indices
$h-2i-1$ in the first case and $2i-h$ in the second case partition
$\ints0{h-1}$, using alternate parities. Thus each nonmultiple of $k$
appears once. Leaf weights are $hk+1,\ldots,hk+m$.
When $h=0$, this is the ordinary star labeling, including a singleton
center when also $m=0$.
\end{proof}

\begin{corollary}\label{cor:leaves}
The center and every existing unit leaf can be assigned zero.
\end{corollary}
\begin{proof}
The center is Lemma~\ref{lem:centerzero}. For a requested leaf, remove it,
label the remainder with center zero, and reattach that exact named leaf
with new label $N$. Its edge contributes the missing weight $N$.
Apply Lemma~\ref{lem:complement} to the entire graph.
\end{proof}

For the selected long-arm tip, retain the other $h=n-1$ arms and all
$m$ unit leaves, use Lemma~\ref{lem:centerzero} on them, and translate
their labels by $R$. Set $Q=hk+m$ and $N=Q+2R$. The selected-arm offsets
are
\[
 T_R=[R,R-1,R-1,R-2,R-2,\ldots,0,0].
\]
Use $T_R[d]$ at even depth $d$ and $N-T_R[d]$ at odd depth.
The center offset is $R$, matching the translated residual center.
The two nonroot offset sides are both $\ints0{R-1}$;
adjacent sums are $2R-1,2R-2,\ldots,0$. Thus new weights are
$Q+1,\ldots,N$, and the terminal vertex has literal zero.
Notice this is a root-$R$ interface, distinct from the root-$R+1$
shells in Section~\ref{sec:shell}.

\section{Auxiliary paths and the center-one residual}
\label{sec:one}
\subsection{A cracked offset path for every radius}
For $r\ge3$ construct a word $A_r$ of length $2r-1$, with even-index
inventory $\{0\}\cup\ints2r$, odd-index inventory $\ints0{r-2}$,
adjacent sums $\ints1{2r-2}$, first entry zero, and last entry $r-1$.
The bases are
\begin{align*}
 A_3&=[0,1,3,0,2],& A_4&=[0,1,2,0,4,2,3],\\
 A_5&=[0,1,3,0,2,3,5,2,4],&
 A_7&=[0,1,3,0,2,3,5,2,4,5,7,4,6].
\end{align*}
For $r=3,5,7$ and thereafter in the corresponding residue class modulo
three, replace the last four entries
$[r-2,r,r-3,r-1]$ by
\[
 [r-2,r-1,r-3,r+1,r-1,r,r+1,r+3,r,r+2].                 \tag{A}
\]
This defines $A_{r+3}$. Radius four is the standalone exception.

\begin{lemma}\label{lem:aux}
All the stated inventories hold, and for every $r\ge5$ the last five
entries are $[r-3,r-2,r,r-3,r-1]$.
\end{lemma}
\begin{proof}
Check the four displayed bases by addition. In (A), the old removed
odd entries are $r-2,r-3$ and the new odd entries are
$r-2,r-3,r-1,r+1,r$; the old removed even entries are $r,r-1$ and
the new ones are $r-1,r+1,r,r+3,r+2$. Thus the sides acquire exactly
the three new values required. The three removed internal sums are
$\ints{2r-4}{2r-2}$, replaced by $\ints{2r-4}{2r+4}$.
The retained prefix join is unchanged. Retained internal sums together
with that join are $\ints1{2r-5}$. At $r=3$ the retained prefix is
just $[0]$: it has no internal edge, but the retained join has sum one.
This boundary contributes the singleton $\{1\}$, not an empty inventory.
The new last four entries have the required terminal form for $r+3$;
the last five follow directly. Bases $3,5,7$ cover all $r\ge3$ except
$4$, which was checked separately.
\end{proof}

\subsection{One and two arms}
For $k=2r$, the outward one-arm path is $[1,k]$ followed by all of
$A_r$, encoding its even-index entries as themselves and odd-index
entries $x$ as $k-1-x$. It has $k+1$ labels, center one, and initial
triple $[1,k,0]$. Its labels are exactly $\ints0k$. The first two
weights are $k-1,k$; the remainder are $k-1-\ES(A_r)$,
namely $\ints1{k-2}$. Hence it is graceful.

For two arms, let
\[
 P_r=\mathop{\cat}_{i=0}^{r-2}[r+1+i,3r-1-i]\cat[2r].
\]
Follow this by $[4r-1,1,4r,0]$, then $A_r$ without its first zero,
encoding odd-index entries $x$ of $A_r$ as $4r-2-x$ and even-index
entries as themselves. The resulting path has $4r+1=2k+1$ vertices,
midpoint label one, and midpoint triple $[1,2k,0]$.
The prefix $P_r$ has weights $\ints1{k-2}$; its next join has weight
$k-1$. The other three exceptional edges have weights $2k-2,2k-1,2k$.
The encoded $A_r$ part has weights $\ints k{2k-3}$.
Labels partition the lower interval $\ints0k$ and the upper interval
$\ints{k+1}{2k}$. Every edge crosses the cut $k$.
Viewing the two halves as the two named arms gives the required
center-one two-arm labeling.

\subsection{Adding long arms}
\begin{lemma}\label{lem:one}
For every even $k\ge6$, $h\ge1$, and $m\ge0$, the actual graph
$S(k^h,1^m)$ has a graceful labeling with center one and maximum label
$hk+m$ at a depth-one vertex whose two neighbors have labels one and zero.
\end{lemma}
\begin{proof}
The one-arm case above establishes $h=1,m=0$. For $h\ge2,m=0$, use
the two-arm base and maintain the following invariant. At cut $T=hk/2$,
the weights on noncrossing edges are exactly
\[
 C_h=\{jk:1\le j\le\lfloor(h-1)/2\rfloor\};
\]
every other edge crosses the cut. The center is one and the maximum
$hk$ remains at depth one with neighbors one and zero.
To add one arm, raise every label greater than $T$ by $k$.
The new labels are the band $T+1,\ldots,T+k$. On the new arm put this
band in Walecki order
$[0,k-1,1,k-2,\ldots]$ translated by $T+1$ if $h$ is even;
reverse that order if $h$ is odd.

The internal new weights are $\ints1{k-1}$ and its center edge has weight
$\lceil h/2\rceil k$. Old crossing weights increase by $k$ and old
noncrossing weights stay fixed. The multiset union
\[
 C_h\ \cup\ ((\ints1{hk}\setminus C_h)+k)
 \ \cup\ \ints1{k-1}\ \cup\{\lceil h/2\rceil k\}
 =\ints1{(h+1)k}
\]
is disjoint, by inspecting the missing multiples of $k$.
At the new cut $T+k/2$, old crossing edges still cross, since their
raised endpoints exceed $T+k$; old noncrossing edges remain on one side.
New internal edges cross. The new root edge is noncrossing exactly when
$h$ is even, giving $C_{h+1}$. The existing maximum is raised by $k$
and keeps its two neighbors. This proves the induction for every $h$.

Now allow $m$ unit leaves. Raise the old maximum $hk$ to $hk+m$ and put
the original named leaves at $hk,hk+1,\ldots,hk+m-1$ on the center.
The two old maximum edges, previously of weights $hk-1,hk$, now have
weights $hk+m-1,hk+m$. The leaf weights are
$hk-1,\ldots,hk+m-2$. The untouched edges fill $\ints1{hk-2}$.
Labels and weights are complete and the maximum anchor survives.
For $m=0$ the operation is empty.
\end{proof}
Depth one on any selected long arm is therefore a zero target: apply
Lemma~\ref{lem:one} to the entire graph, permute the anchored arm to
the requested name, and complement every label.

\section{Balanced shells and their actual graph transfer}
\label{sec:shell}
\begin{definition}
A shell of integer radius $q\ge3$ is a word $O=(O_0,\ldots,O_{2q})$ with
$O_0=q+1$, each of its nonroot odd- and even-index inventories equal to
$\ints0{q-1}$, and $\ES(O)\mult\ints0{2q-1}$.
\end{definition}
\begin{lemma}[Forced boundaries]\label{lem:forced}
Every shell has $O_1=q-2$ and $O_{2q}=1$.
\end{lemma}
\begin{proof}
Internal nonroot sums are at most $2q-2$, so sum $2q-1$ must be the
root edge $(q+1)+O_1$. For the endpoint, sum all adjacent sums:
\[
 q(2q-1)=2\bigl(q(q-1)+q+1\bigr)-(q+1)-O_{2q},
\]
which gives $O_{2q}=1$.
\end{proof}

\begin{lemma}[Shell transfer]\label{lem:transfer}
Let $k=2R$ with $R\ge3$, $n\ge2$, $m\ge0$, and let $O$ be a shell of radius $R$.
If $O_j=0$ at a nonroot index $j$, then every selected arm has a graceful
labeling putting zero at depth $j$.
\end{lemma}
\begin{proof}
Use Lemma~\ref{lem:one} on the other $h=n-1\ge1$ named arms and all
original unit leaves. Put $Q=hk+m$, $N=Q+k$ and translate these residual
labels by $R$. Their center is $R+1$, equal to $O_0$.
The selected-arm labels are $O_d$ at even depth and $N-O_d$ at odd depth.
The three label intervals are
\[
 \ints0{R-1},\qquad\ints R{Q+R},\qquad\ints{Q+R+1}N.
\]
Each new edge has positive weight $N-(O_{d-1}+O_d)$, so the new weights
are $\ints{Q+1}N$ and the unchanged residual weights are $\ints1Q$.
If $j$ is even its zero is literal; if odd its label is $N$, and the
whole graph is complemented. The graph names, all arm lengths, and all
center-leaf edges are unchanged. The $h=1$ anchored path handles
$n=2$ for every $m$, including $m=0$.
\end{proof}

% V2_CONTEXT_BEGIN attachment_credit

The attachment mechanism in Lemma~\ref{lem:transfer} agrees with
Shan--Zhong's endpoint-path attachment~\cite[Theorem 2]{sz} at a path
on $L=k$ vertices and residual root label one: its condition is
$1+\lfloor L/2\rfloor+1\le L$ and $L\not\equiv1\pmod4$.
Their shifts match the label intervals above. The shell supplies the
additional simultaneous endpoint and prescribed internal extreme
needed for the selected zero; ordinary path attachment alone does not
choose that position.
% V2_CONTEXT_END attachment_credit

For $r\ge4$, apply Lemma~\ref{lem:aux} to $A_{r+1}$ and encode its
even positions by $r+1-x$ and odd positions by $r-1-x$.
This yields a canonical shell $B_r$ with terminal five entries
$[3,0,0,1,1]$. Indeed each side is $\ints0{r-1}$, the root is $r+1$,
and each sum is $2r-\ES(A_{r+1})$, giving $\ints0{2r-1}$.

\begin{lemma}[Graft]\label{lem:graft}
Let $O$ have radius $q$ and let $r\ge4$. Delete the last four entries
of $B_r$, shift the retained entries by $q-2$, and append $O$ without
its root. The resulting word $G$ is a shell of radius $R=r+q-2$.
An entry at source index $j\ge1$ moves to index $j+2r-4$.
Its deficit $2q-j$ is preserved.
\end{lemma}
\begin{proof}
The retained word has length $2r-3$, root $R+1$, and last entry $q+1$.
On each nonroot parity the retained offsets are $\ints q{R-1}$;
the suffix supplies $\ints0{q-1}$. Removing the four terminal edges
removes sums $3,0,1,2$ from $B_r$. After translation, retained internal
sums become $\ints{2q}{2R-1}$. The join is
$(q+1)+O_1=2q-1$ by Lemma~\ref{lem:forced}.
The root-deleted source sums are $\ints0{2q-2}$. All three bands are
disjoint and complete. The length gives the stated position shift;
$2R-(j+2r-4)=2q-j$.
\end{proof}
The statement applies to any valid shell because its first high is
forced. A contract giving only the internal suffix sums would not justify
that join.

After deleting the common root and using the parity encoding, this
operation is the translated balanced concatenation in the proof of
Hicks--Ollis--Schmitt's Lemma 4.5~\cite{hos}. Its compatible prefix
endpoints are also available from Ollis's Lemma 5.5~\cite{ollis},
credited there to Kotzig, reference~[21]. The formula above checks the additional
source-index transport explicitly. The graft operation itself is not
claimed as new; Ollis's Theorem 5.6 is a separate insertion construction,
with its own stated adjacent-entry hypotheses.

\section{Balanced midpoint cores and the prefix interface}
\label{sec:cores}
\begin{definition}
A boundary core of size $p\ge4$ is a length-$2p$ word $C$ whose
even-index and odd-index entries each permute $\ints0{p-1}$,
whose sums are $\ints0{2p-2}$, and whose first and last entries are
$3,p-4$. Its even indices are called high and its odd indices low.
An anchored core has $[4,0,0,1]$ at indices $d-2,d-1,d,d+1$, for some
even $d$; in particular its two zeros have outward depths $d,d+1$.
\end{definition}

\begin{lemma}[Midpoint alpha lift]\label{lem:corelift}
A boundary core of size $p$ yields zero at every outward depth $d$ with
$C_{d-1}=0$ on any selected arm of the actual graph
$S((2p+4)^n,1^m)$, for all $n\ge2,m\ge0$.
\end{lemma}
\begin{proof}
Set $k=2p+4$ and $M=2k$. Define two outward arms of length $k$:
\begin{align*}
 S&=(M-C_0,C_1,M-C_2,C_3,\ldots,C_{2p-1})
       \cat[3p+5,p+2,3p+7,p+3],\\
 T&=\mathop{\cat}_{i=0}^{p-1}[2p+5+i,2p+3-i]
       \cat[3p+6,p,3p+8,p+1].
\end{align*}
Put the common center at $k$. The full path is $S$ reversed, then $k$,
then $T$. The label inventory is recorded here by side:
\begin{center}\begin{tabular}{c ll}\toprule
 & lower labels & upper labels\\\midrule
 $S$ & $\ints0{p-1},p+2,p+3$ & $\ints{3p+9}{4p+8},3p+5,3p+7$\\
 $T$ & $p,p+1,\ints{p+4}{2p+3}$ & $\ints{2p+5}{3p+4},3p+6,3p+8$\\
 center & $2p+4$ & \\\bottomrule\end{tabular}\end{center}
Core weights are $\ints{2p+10}{4p+8}$; the selected center edge and
four tail edges have weights $2p+1$ and $2p+9,2p+3,2p+5,2p+4$.
The partner center edge and prefix internal edges give $\ints1{2p}$;
its four tail edges give $2p+2,2p+6,2p+8,2p+7$. Thus all weights are
$\ints1M$ and every edge crosses the cut $k$.

Now put $h=n-2$, $Q=hk+m$, $N=M+Q$. Translate the
center-zero residual from Lemma~\ref{lem:centerzero} by $k$ and identify
its designated center with this midpoint. Leave alpha-path labels
$\le k$ fixed and raise labels $>k$ by $Q$. Its weights increase by
$Q$ because every edge crosses the cut; the residual contributes
$\ints1Q$. Labels partition $\ints0{k-1}$, $\ints k{k+Q}$,
$\ints{k+Q+1}N$. A low core zero is literal graph zero; a high core
zero is $N$ and uses full complementation. When $h=m=0$ the residual
is just the center; when $h=0,m>0$ it is the original center star.
No leaf has been subdivided or moved.
\end{proof}

% V2_CONTEXT_BEGIN amalgamation_credit

The final residual splice in Lemma~\ref{lem:corelift} is classical
zero/index amalgamation of Huang--Kotzig--Rosa, as precisely restated
in Shan--Zhong~\cite[Lemma 1]{sz}. The original 1982 full proof was
not obtained here. The explicit midpoint alpha witness above supplies
the joint center-at-index and internal-zero conditions; the
amalgamation operation itself is not presented as new.
% V2_CONTEXT_END amalgamation_credit

\section{A self-contained complete prefix}
\label{sec:prefix}
This section proves the older compatible prefix $P$ directly. It is
sufficient for the main theorem and is kept distinct from the stronger
current $D_8$ in Section~\ref{sec:d8}.

\subsection{Finite seeds and size-six extension}
Appendix~\ref{app:cores} displays all $30$ boundary cores
$B_{p,d}$ for $8\le p\le13$ and $d\in\{2,4,6,8,10\}$, with zeros
at indices $d-1,d$. It also displays all $48$ anchored cores $C_{p,d}$
for $16\le p\le21$, $d\in\{12,14,16,18,20,22,24,26\}$.
These finite certificates satisfy exactly the preceding definitions by
parity inspection and adjacent addition; no search claim is required.
Put
\[
 C_6=[3,3,2,5,5,4,4,0,0,1,1,2].
\]
If $C$ is a boundary core of size $p$, then
$C\cat(p+C_6)$ is a boundary core of size $p+6$.
The two side inventories acquire $\ints p{p+5}$.
The old sums are $\ints0{2p-2}$, the join is
$(p-4)+(p+3)=2p-1$, and shifted new sums are $\ints{2p}{2p+10}$.
Existing zero positions and an anchored window stay unchanged.
For every $p\ge8$, select $p_0=8+(p-8)\bmod6$ and extend
$B_{p_0,d}$ repeatedly. Hence every depth $2,\ldots,11$ is available
for $k=2p+4$.

\subsection{Insertion cards and their exact inventories}
Let an anchored core be $C=u\cat[0,0]\cat v$, where $|u|=d-1$ is
odd, $u$ begins at $3$ and ends at $4$, and $v$ begins at $1$ and ends
at $p-4$. All entries of $u,v$ are positive. An insertion card is
$(q;P,B,D)$ with all three lengths even, total length $2q$,
$P_0=3$, last entry of $B$ equal to $4$, and first entry of $D$ equal to
$1$. The inserted word is
\[
 P\cat(q+u)\cat B\cat[0,0]\cat D\cat(q+v).             \tag{I}
\]
Its zero depths increase by $\delta=|P|+|B|$.
The literal cards are in Appendix~\ref{app:cards}.
For each card their fresh high inventory
$\Hi(P)\cat\Lo(B)\cat\Lo(D)$ and fresh low inventory
$\Lo(P)\cat\Hi(B)\cat\Hi(D)$ are both $\ints1q$.
Here $\Hi$ extracts even positions in the indicated isolated word;
$B,D$ start on low parity in (I), explaining the swaps.
The three fresh edge chains have inventory
\[
 \ES(P\cat[q+3])\cat
 \ES([q+4]\cat B\cat[0])\cat
 \ES([0]\cat D\cat[q+1])
 \mult\ints1{2q+1}\cup\{2q+4\}.                       \tag{J}
\]

\begin{lemma}\label{lem:insert}
Each displayed card transforms an anchored core of size $p$ into one of
size $p+q$, preserving its anchored window and increasing its zero depths
by $\delta$.
\end{lemma}
\begin{proof}
The two original zeros remain. Original positive offsets $1,\ldots,p-1$
on each side shift to $q+1,\ldots,p+q-1$, and fresh values supply
$1,\ldots,q$. The endpoint and window statements follow from the card
boundaries. Removed old sums are $4,0,1$; the zero sum is retained.
The untouched sums shift by $2q$, giving
$\{2q+2,2q+3\}\cup\ints{2q+5}{2(p+q)-2}$.
Together with zero and (J), these are exactly
$\ints0{2(p+q)-2}$. All bands are disjoint. The new low-zero index is
$|P|+|u|+|B|$, proving the depth increase.
\end{proof}
Six cards have $q=9$ and $\delta=4,6,8,10,12,14$; the seventh has
$q=18$ and $\delta=30$. Two size-nine cards have every even
displacement $8,10,\ldots,28$, so with the size-eighteen card one macro
step of size eighteen has all displacements $8+2j$, $0\le j\le11$.
Write $t=2a+e$, $e\in\{0,1\}$. Using $a$ macro steps and $e$
size-nine steps realizes every even displacement
\[
 4t+2j,\qquad 0\le j\le11a+5e=5t+\lfloor t/2\rfloor.  \tag{K}
\]
This is the sum of integer intervals, not an extrapolation from seeds.

\begin{proposition}[Complete prefix]\label{prop:prefix}
For every even $k\ge20$, every $n\ge2,m\ge0$, every selected arm,
and every depth $2\le d\le P(k)$, zero is available, where
\[
 P(k)=\begin{cases}
 11,&20\le k\le34,\\
 27+14u+2\lfloor u/2\rfloor,&k\ge36,
 \quad u=\lfloor(k-35)/18\rfloor.
 \end{cases}                                           \tag{P}
\]
Moreover $k/2\le P(k)<k$ for every even $k\ge36$.
\end{proposition}
\begin{proof}
The small-depth construction already covers $2,\ldots,11$.
Put $p=(k-4)/2$ and $u=\lfloor(p-16)/9\rfloor$ for $k\ge36$.
Starting with even depths $12,14,\ldots,26$ and their neighboring odd
depths, (K) gives the complete interval
$\ints{12+4t}{27+14t+2\lfloor t/2\rfloor}$ at added size $9t$.
These intervals overlap consecutively: for $t\ge1$, the preceding
upper endpoint is at least $12+4t-1$. Their union for $0\le t\le u$
is $\ints{12}{P(k)}$.

For clarity an exact size choice is necessary. Given a requested $d\ge12$,
choose the least $t$ whose upper endpoint reaches $d$.
Then $t\le u$ and $d\ge12+4t$. Set
$w=\lfloor(d-12-4t)/2\rfloor$, $b=\min(w,7)$, and $j=w-b$.
The starting depth $12+2b$ lies between $12$ and $26$, and
$0\le j\le5t+\lfloor t/2\rfloor$. Choose
\[
 p_0=16+((p-16-9t)\bmod6)\in\ints{16}{21}.
\]
Since $p\ge16+9t$, the remaining size $p-p_0-9t$ is a nonnegative
multiple of six. Start with $C_{p_0,12+2b}$, apply the insertion steps
realizing $j$ in (K), and append $C_6$ for this remaining size.
The resulting core has exactly size $p$; its zero pair covers $d$,
with parity selecting the appropriate member. Lemma~\ref{lem:corelift}
gives the full graph claim for all $n,m$.
Finally even $k$ with the given $u$ satisfies $k\le52+18u$. Thus
$k/2\le26+9u<27+14u+2\lfloor u/2\rfloor=P(k)$.
Also $k\ge36+18u$ gives
$k-P(k)\ge9+4u-2\lfloor u/2\rfloor>0$; hence every listed depth is
inside the actual arm. This establishes the unbounded high branch algebraically.
\end{proof}

\section{Two infinite shell source families}
\label{sec:families}
\subsection{Ordinary and signed caps}
For each $h\ge3$ let $H_h$ be a shell of radius $h$ defined from the
two bases $H_3=[4,1,0,0,2,2,1]$ and
$H_4=B_4=[5,2,2,3,3,0,0,1,1]$ by
\[
 H_{h+2}=[h+3,h,h,h+1]\cat H_h.                       \tag{C}
\]
The old root becomes a new even nonroot entry. Each side gains $h,h+1$;
the four new sums are $2h+3,2h,2h+1,2h+2$. Therefore (C) preserves
the complete shell contract for every $h\ge3$ by the two parity bases.

The signed cap $E_h$ instead has length $2h+3$, root $h+1$, odd-side
inventory $\{-1\}\cup\ints0{h-1}$, even-side inventory
$\{-2\}\cup\ints0{h-1}$, sums $\ints{-2}{2h-1}$, and terminal $-2$.
Use
\[
 E_3=[4,1,1,2,2,-1,0,0,-2],\quad
 E_4=[5,2,3,3,1,-1,0,1,2,0,-2],
\]
and exactly (C) with $E$ in place of $H$. The same four new sums
prove the signed contract for every $h\ge3$. The negative exceptional
entries stay on their own sides. In each cap, the first high $h-2$ is
forced by its maximal sum. Only this minimum signed contract is used;
no terminal-four motif is needed.

\subsection{Valleys and exact side/sum partitions}
Let $t\ge5$, $h=t-2$, $v=3t+2$, and define
\begin{align*}
 L_t&=\mathop{\cat}_{s=t}^{0}[3s+1,3s],&
 A_t&=\mathop{\cat}_{s=0}^{t-1}[3s+2,3s+3],\\
 D_t&=[3t+1]\cat\mathop{\cat}_{s=t-1}^{0}[3s+2,3s+1].
\end{align*}
Decreasing upper-to-lower indices indicate decreasing concatenation.
The first entry of $D_t$ is essential. Put
\begin{align*}
 V_t^+&=L_t\cat[0]\cat A_t\cat[v,v]\cat D_t,\\
 V_t^-&=L_t\cat[0]\cat A_t\cat[v,v+1,v+2,v]\cat D_t.
\end{align*}
The valley begins on high parity in the complete shell. The contributions
from $L_t$ are high $3s+1$, low $3s$; the extra zero is high;
$A_t$ contributes high $3s+3$, low $3s+2$; the descending pairs
in $D_t$ contribute high $3s+2$, low $3s+1$.
The ordinary plateau and $D$ head contribute high $3t+2$, low
$3t+2,3t+1$. Hence each valley side in $V_t^+$ is $\ints0{3t+2}$.
For $V_t^-$ the corresponding contributions are high $3t+3,3t+2$,
low $3t+2,3t+4,3t+1$. Thus its high side is $\ints0{3t+3}$,
and its low side is $\ints0{3t+2}\cup\{3t+4\}$.

The ordinary internal valley sums form $\ints0{6t+4}$. One explicit
check by residue modulo six is: residues $0,1,2,3,4$ give $6s+j$ for
$0\le s\le t$, while residue five gives $6s+5$ for $0\le s<t$.
Here $L$ within-pairs give residue one, $L$ inter-pairs residue four,
$A$ within-pairs residue five and $A$ inter-pairs residue two,
$D$ within-pairs residue three and $D$ inter-pairs residue zero.
The zero and plateau joins fill the respective boundary members.

For the signed valley, include the cap-to-valley join $6t+4$.
The resulting sums are $\ints0{6t+7}$, partitioned as follows:
\begin{center}\begin{tabular}{c c l}\toprule
 residue & allowed $s$ in $6s+j$ & additional boundary contributors\\\midrule
0 & $0\le s\le t+1$ & zero, $D$ bridge $6t$, plateau $6t+6$\\
1 & $0\le s\le t+1$ & plateau $6t+7$\\
2 & $0\le s\le t$ & zero-to-$A$, $A$-to-plateau\\
3 & $0\le s\le t$ & plateau $6t+3$\\
4 & $0\le s\le t$ & cap-to-$L$ join $6t+4$\\
5 & $0\le s\le t$ & plateau $6t+5$\\\bottomrule
\end{tabular}\end{center}
The pair contributors are as in the ordinary case. The table lists every
edge once, including the exceptional joins; it is a multiset partition.

\begin{theorem}[Uniform odd-radius sources]\label{thm:source}
Every odd $q\ge21$ has a shell $U_q$ whose zeros are at indices
$q-3$ and $q-2$.
\end{theorem}
\begin{proof}
For $q=4t+1$ set $a=3t+3$ and
\[
 U_q=[q+1]\cat(a+H_h[1:])\cat V_t^+.
\]
The shifted cap gives each nonroot side $\ints{3t+3}{4t}$; the valley
gives $\ints0{3t+2}$. Shifted cap sums are
$\ints{6t+6}{8t+1}$, the cap-valley join is $6t+5$, and valley sums
are $\ints0{6t+4}$. The full inventory is $\ints0{2q-1}$.
The length is $(2t-3)+(6t+6)=8t+3=2q+1$.
The low zero at the end of $L_t$ has index $4t-2=q-3$ and the next
high zero has index $4t-1=q-2$.

For $q=4t+3$ set $a=3t+5$ and
\[
 U_q=[q+1]\cat(a+E_h[1:])\cat V_t^-.
\]
The cap high inventory is $\ints{3t+4}{4t+2}$ and its low inventory
is $\{3t+3\}\cup\ints{3t+5}{4t+2}$. The valley inventories above
complete both sides to $\ints0{q-1}$. Shifted signed cap sums are
$\ints{6t+8}{8t+5}$; the remaining sums including the join are
$\ints0{6t+7}$. Thus full sums are $\ints0{2q-1}$.
The length is $(2t-1)+(6t+8)=8t+7=2q+1$; zeros have indices
$4t=q-3$ and $4t+1=q-2$. The cap shift includes its root because
$a+h+1=q+1$ in both cases. Every odd $q\ge21$ is exactly one of
these two forms with $t\ge5$.
\end{proof}

\begin{corollary}[All large deficits]\label{cor:deficits}
For an even ambient length $k=2R$, a requested deficit $c=k-d\ge23$
is available whenever $c\le R-1$.
\end{corollary}
\begin{proof}
For odd $c$, use $q=c-2$ and source high-zero index $j=q-2$;
for even $c$, use $q=c-3$ and low-zero index $j=q-3$.
Then $q$ is odd and at least $21$, and $q\le R-3$.
Set $r=R-q+2\ge5$. The graft index is
$j+2r-4=2R-c=d$. Apply Lemma~\ref{lem:transfer}, with full
complementation in the high case.
\end{proof}

\section{Finite tail sources and lengths below thirty-six}
\label{sec:finite}
\subsection{The last eight depths}
Appendix~\ref{app:shells} lists near-tip seeds of radii $5,6,9$.
Each is a shell. Prefixing $[r+3,r,r,r+1]$ increases its radius by two
and every old nonroot index by four, preserving deficit.
The radius-six seeds supply pairs of deficits $(1,2),(3,4),(5,6),(7,8)$
for all even radii $R\ge8$. For odd radii $R\ge9$, the radius-five
seeds supply $(1,2),(3,4),(4,5),(6,7)$ and the radius-nine seed supplies
$(7,8)$. Their union supplies every deficit $1,\ldots,8$.
Apply Lemma~\ref{lem:transfer}; no claim of odd-depth impossibility is
suggested by the use of low zeros on the even depths.

\subsection{Deficits nine through twenty-two}
The source bank in Appendix~\ref{app:shells} has the following exact
zeros and useful deficits. An entry named $q\mathrm a$ or $q\mathrm b$
has radius $q$; the letters distinguish witnesses only.
\begin{center}\begin{tabular}{c c c c}\toprule
 source & zero indices & deficits & sufficient ambient $R$\\\midrule
 $9\mathrm b$ & $8,9$ & $10,9$ & $R\ge11$\\
 $9\mathrm a$ & $6,7$ & $12,11$ & $R\ge11$\\
 $11\mathrm b$ & $9,10$ & $13,12$ & $R\ge13$\\
 $11\mathrm a$ & $7,8$ & $15,14$ & $R\ge13$\\
 $12$ & $8,9$ & $16,15$ & $R\ge14$\\
 $15$ & $12,13$ & $18,17$ & $R\ge17$\\
 $17$ & $14,15$ & $20,19$ & $R\ge19$\\
 $21$ & $20,21$ & $22,21$ & $R\ge23$\\\bottomrule
\end{tabular}\end{center}
These bounds use Lemma~\ref{lem:graft}. A source can also be used
directly at $R=q$, without asserting a canonical radius-two or
radius-three graft. Several deficits have alternative valid witnesses.

\subsection{Explicit small-length coverage}
Small extra midpoint cores in Appendix~\ref{app:cores} cover depths
$2$--$7$ at $k=18$, and $12$--$15$ at $k=32,34$.
For $k=20,22,24,26,28,30,32,34$, Proposition~\ref{prop:prefix}
supplies depths $2$--$11$. The following table lists exact additional
source placements; ``direct'' means using a shell at its own radius,
and $r$ means its canonical graft radius.
\begin{center}\begin{tabular}{c p{105mm}}\toprule
$k$ & additional coverage before the last-eight interval\\\midrule
18 & extra $p=7$ cores: $2$--$7$; $9\mathrm b$ direct: $8,9$.\\
20 & no additional depths needed.\\
22 & $9\mathrm b$, $r=4$: $12,13$.\\
24 & $9\mathrm a$, $r=5$: $12,13$; $9\mathrm b$, $r=5$: $14,15$.\\
26 & $11\mathrm a$, $r=4$: $11,12$; $11\mathrm b$, $r=4$: $13,14$;
      $9\mathrm a$, $r=6$: $14,15$; $9\mathrm b$, $r=6$: $16,17$.\\
28 & $12$, $r=4$: $12,13$; $11\mathrm a$, $r=5$: $13,14$;
      $11\mathrm b$, $r=5$: $15,16$; $9\mathrm a$, $r=7$: $16,17$;
      $9\mathrm b$, $r=7$: $18,19$.\\
30 & $15$ direct: $12,13$; $12$, $r=5$: $14,15$;
      $11\mathrm a$, $r=6$: $15,16$; $11\mathrm b$, $r=6$: $17,18$;
      $9\mathrm a$, $r=8$: $18,19$; $9\mathrm b$, $r=8$: $20,21$.\\
32 & extra $p=14$ cores: $12$--$15$; $12$, $r=6$: $16,17$;
      $11\mathrm a$, $r=7$: $17,18$; $11\mathrm b$, $r=7$: $19,20$;
      $9\mathrm a$, $r=9$: $20,21$; $9\mathrm b$, $r=9$: $22,23$.\\
34 & extra $p=15$ cores: $12$--$15$; $15$, $r=4$: $16,17$;
      $12$, $r=7$: $18,19$; $11\mathrm a$, $r=8$: $19,20$;
      $11\mathrm b$, $r=8$: $21,22$; $9\mathrm a$, $r=10$: $22,23$;
      $9\mathrm b$, $r=10$: $24,25$.\\\bottomrule
\end{tabular}\end{center}
All entries follow the single position formula $j\mapsto j+2r-4$.
Their unions, with depths $k-8,\ldots,k-1$, cover every depth $2$ through
$k-1$. The named center, leaves, depth one and tip use the separate
constructions above. This proves every fixed even length $18$--$34$
without treating a model-specific failed shell search as graph impossibility.

\section{Completion of the unbounded theorem and the D8 seam}
\label{sec:d8}
\begin{proof}[The $k\ge18$ branch of Theorem~\ref{thm:main}]
Small lengths were proved in Section~\ref{sec:finite}. Let even $k\ge36$.
For a requested long-arm depth $d$, handle $d=1,k$ by the endpoint
constructions, and $2\le d\le P(k)$ by Proposition~\ref{prop:prefix}.
For the remaining depths write $c=k-d$. The midpoint inequality gives
$1\le c\le k/2-1$. If $c\ge23$, Corollary~\ref{cor:deficits} applies.
For $1\le c\le22$ use the finite tail bank. At $k=36,38,40,42,44$
the largest required deficit is $8,10,12,14,16$, respectively, so all
needed source bounds fit. At $k=46,48,50$ the largest required deficits
are $18,20,22$ and the relevant new source radii are $15,17,21$.
At $k=52$ the largest deficit is $24$ and $q=21$ fits. For $k\ge54$
all finite source bounds hold. The shell transfer gives every actual
arm vertex, and Corollary~\ref{cor:leaves} gives the center and leaves.
Thus every named vertex is covered for every $n\ge2,m\ge0$.
\end{proof}

\subsection{Relation to the stronger recorded prefix}
The prior checked prefix chain denoted $D_8$ has the following current
even-domain values:
\[
 D_8(k)=\begin{cases}
 11,&20\le k\le34,\\
 27,&36\le k\le52,\\
 25+16\lfloor(k-35)/18\rfloor,&54\le k\le124,\\
 107+20a+2s,&k=126+22a+2s,\quad a\ge0,\ 0\le s\le10.
 \end{cases}                                           \tag{D8}
\]
Its domain in the recorded construction is $k\ge19$, but (D8) is
written only for even lengths here. At $18$ it is not invoked.
The current chain has extra insertion/reflection gap fills beyond (P).
Neither (D8) nor (P) is invoked for the fixed lengths $6,8,10,12,14,16$;
the certified words and transfers in Section~\ref{sec:below18} handle them.
We do not identify (P) with (D8). For example, at $352$ they are
$281$ and $313$, respectively.

The integrated proof using $D_8$ has the same tail seam:
$D_8(k)\ge k/2$ for all even $k\ge36$. This follows directly from
(D8): the middle blocks have values $41,57,73,89$ and maximal half-lengths
$35,44,53,62$; the high branch satisfies
$D_8(k)-k/2=44+9a+s$. Also $D_8\ge P$ throughout their even domain.
For $54\le k\le124$, the difference is
$-2+2u-2\lfloor u/2\rfloor\ge0$; the first two branches coincide.
In the high branch, $P(k)\le(5k-13)/6$ while
$6D_8(k)-5k+13=25+10a+2s>0$.
The self-contained proof above uses the weaker (P), so it does not require
readers to reconstruct the extra $D_8$ modules. These formulas describe
the exact relation to the previously integrated seam and preserve its
stronger, separate scope.

\section{The complete small cases two and four}\label{sec:twofour}
These cases use their own constructions. Neither specializes the
radius-$R\ge3$ shell transfer, the $r\ge4$ canonical graft, or any
large-length prefix below its stated domain. All arm and original leaf
names remain those of Section~\ref{sec:statement}.

\subsection{Length two: every named target and every leaf count}
For $S(2^h,1^s)$ with $h\ge1$, Lemma~\ref{lem:centerzero} is
\[
 b(c)=0,\quad b(v_{i,1})=2(h-i),\quad b(v_{i,2})=2i+1,
 \quad b(u_j)=2h+j+1.
\]
The root edges give all positive even weights through $2h$.
The outer weights are $|2h-4i-1|$: they are positive odd numbers at
most $2h-1$, and equality with opposite signs would imply
$4(i+j)=4h-2$, impossible modulo four. Thus all labels and weights
are complete for every $h,s$, including the one-arm residual.

Put $N=2n+m$. For requested support $v_{a,1}$, delete exactly its
named two-edge arm, retain every original unit leaf, and use this
source on $h=n-1\ge1$ arms, with $N-2$ edges. Relabel every retained
vertex by $N-1-b(x)$, give $v_{a,1}$ label zero and $v_{a,2}$ label
$N$. The retained labels are $\ints1{N-1}$, retained weights
$\ints1{N-2}$, and center label $N-1$ gives the two missing weights
$N-1,N$. Whole-graph complement puts zero at the selected tip
$v_{a,2}$. The center uses the original source.

For a requested original leaf $u_j$, necessarily $m\ge1$, delete
only that named leaf and use the center-zero source on the remainder,
of size $N-1$. Reflect its labels by $N-b(x)$ and give the deleted
leaf zero: retained labels fill $\ints1N$, retained weights
$\ints1{N-1}$, and the new center-to-leaf difference is $N$.
At $m=0$ this target class is empty. Other retained arm and leaf
indices may be compressed in increasing order for evaluating the
source, then restored to their original names.

For $n=2,m=0$, a selected-support labeling in the order
$(c,v_{0,1},v_{0,2},v_{1,1},v_{1,2})$ is $(3,0,4,1,2)$.
For $n=2,m=1$ it is $(4,0,5,2,3,1)$ after adjoining the original
leaf; zero at that original leaf is supplied by $(5,1,4,3,2,0)$.
These boundary illustrations do not replace the universal interval proof.
The symmetric $m=0$ case and the $n=2,m=0/1$ cases already have
earlier coverage~\cite{rofa,lcr}. The broader Van Bussel comparison
depends on an inaccessible original exception-class text; no definitive
attribution or novelty inference is drawn from that access gap.

\subsection{Length four: a complete deduction from older operations}
We give the sufficient old-operation composition for every target,
including all $m\ge2$. This does not assert that an earlier paper
literally stated this entire all-$n,m$ corollary.

\paragraph{Center and original leaves.}
Krop's account of the classical generalized $\Delta_{+1}$
construction~\cite[Theorem 2.4 and subsequent generalization]{krop}
starts from a star with exceptional hub label $h$ and leaf labels
$i=0,\ldots,h-1$. Substitute a four-vertex Walecki path
$(0,3,1,2)$ rooted at its zero endpoint into each nonexceptional
vertex. The copy labels are
\[
 (4i,4(h-1-i)+3,4i+1,4(h-1-i)+2),\qquad c=4h.
\]
The hub-to-copy-root edge followed by its three path edges is
exactly a four-edge arm. Complementing by $4h$ gives hub zero and
arm $i$ labels $(4(h-i),4i+1,4(h-i)-1,4i+2)$, the $k=4$ instance
of Lemma~\ref{lem:centerzero}. Leaves $4h+1,\ldots,4h+m$ extend
this hub-zero graph gracefully. Removing and reattaching the exact
requested leaf, then complementing, gives its zero by
Corollary~\ref{cor:leaves}.

\paragraph{Depths one and two, with all original leaves.}
Start with the center-zero $S(2^h)$ labels
$(2(h-i),2i+1)$ on arm $i$. Select its maximum $2h$ at arm-zero
depth one as exceptional vertex in a second two-vertex-copy
$\Delta_{+1}$ construction. After complementing, the old vertex
$x$ other than that exception has copy labels
\[
 A_x=4h-2b(x),\qquad B_x=2b(x)+1,
\]
and the exception has label zero. Identify $B_c$ as the new actual
hub, label one. The selected arm is
$(A_c,\text{exception},A_{0,2},B_{0,2})$, with labels
$(4h,0,4h-2,3)$. Other arms $i>0$ have labels
\[
 (4(h-i)+1,4i,4(h-i)-2,4i+3).
\]
The exception connects only to the prescribed copy roots $A_c$ and
$A_{0,2}$. Other intercopy joins are $B_cB_{i,1}$ and
$A_{i,1}A_{i,2}$, corresponding vertices in the respective copies.
Thus the only branching hub is $B_c$, and each actual arm has four
edges; the generalized construction's attachment conditions are met.

The maximum $4h$ is now the selected depth-one vertex with neighbors
one and zero. Raise just this maximum to $4h+m$ and assign the
original hub leaves $4h,\ldots,4h+m-1$. Untouched weights are
$\ints1{4h-2}$; leaf weights are $4h-1,\ldots,4h+m-2$; the two
maximum-edge weights become $4h+m-1,4h+m$. This is a disjoint
partition, and the vertex labels likewise remain complete. For $m=0$
the operation is empty. Set $h=n$ and permute equal-arm names to
the specified arm. Depth two remains zero; whole-graph complement
gives depth one zero. This explicitly proves the unrestricted
leaf extension rather than inferring it from the symmetric $m=0$ case.

\paragraph{Depths three and four.}
Delete the requested four-edge arm. The actual residual
$S(4^{n-1},1^m)$ has a hub-zero labeling $b$ and
$Q=4(n-1)+m$ edges. The path $(2,3,1,4,0)$ is an alpha labeling
of $P_5$ with cut two and joining endpoint two. It is the inverse
alpha labeling of the Walecki path $(0,4,1,3,2)$.
Identify that endpoint with the residual hub, shift all residual
labels by two, and raise the high path labels by $Q$. This is the
classical zero/index identification of Huang--Kotzig--Rosa, precisely
restated in Shan--Zhong~\cite[Lemma 1]{sz}; no direct original 1982
proof inspection is claimed. With $N=Q+4$, the selected arm labels
are $(N-1,1,N,0)$ and the actual hub is two. Residual labels fill
$\ints2{Q+2}$; the four added labels are $0,1,N-1,N$.
Residual weights stay $\ints1Q$, and new weights in outward order
are $Q+1,Q+2,Q+3,Q+4$. This gives depth four zero and depth three
maximum, hence depth three zero after full complement. Every original
unit leaf stays incident to the actual shared hub.

The center, each existing named leaf, and depths one through four
on each named arm are now covered. The proofs include $n=2,m=0/1$
and allow every natural $m$. Neither an invalid radius-two shell
specialization nor a size-one prefix graft is used. The full $k=4$
construction is a deduction from older sufficient operations; it is
not claimed as an original contribution. Earliest literal statement
and priority for the complete even-length union remain unknown.

\section{The fixed even lengths six through sixteen}\label{sec:below18}
This section supplies the six lengths outside the preceding $k\ge18$
argument. None uses $P(k)$, whose stated proof begins at $20$, or $D_8$,
whose recorded domain begins at $19$. The eight-depth shell theorem has
radius at least eight; below $k=16$ we check each smaller radius directly.
All words below are finite, literal certificates. Their listed side and
sum inventories can be checked without a search algorithm.

\paragraph{Two transfer rules and the two-arm boundary.}
First, if a path word $W_0,\ldots,W_{2k}$ permutes $\ints0{2k}$, has
$W_k=k$, all edges crossing cut $k$, and edge weights $\ints1{2k}$,
choose a named selected arm as the reversed left half and another arm as
the right half. Set $Q=k(n-2)+m$ and $N=2k+Q$. Translate the center-zero
residual $S(k^{n-2},1^m)$ by $k$, preserve the path labels at most $k$,
and raise its labels above $k$ by $Q$. The label bands are
$\ints0{k-1}$, $\ints k{k+Q}$, $\ints{k+Q+1}N$; residual weights are
$\ints1Q$ and path weights become $\ints{Q+1}N$. A path zero on the
selected half remains zero; a path maximum becomes $N$ and is made zero
by complementing \emph{the whole graph}. At $n=2,m=0$ the residual is
only the specified midpoint ($Q=0$). At $n=2,m=1$ it is the midpoint
with the original unit leaf ($Q=1$). Both are literal cases of the
same interval argument.

Second, a radius-$R$ shell $O_0,\ldots,O_{2R}$ has $O_0=R+1$,
both nonroot parity subsequences permuting $\ints0{R-1}$, and adjacent
sums permuting $\ints0{2R-1}$. Lemma~\ref{lem:transfer} applies for
$R\ge3$. Put $Q=2R(n-1)+m$, $N=Q+2R$, translate the center-one residual
on the other $n-1$ named arms and all original leaves by $R$, and put
$O_j$ at even selected depth $j$, $N-O_j$ at odd depth. Its labels
fill $\ints0{R-1}$ and $\ints{Q+R+1}N$ outside the translated residual
$\ints R{Q+R}$; its new weights are
$N-(O_{j-1}+O_j)$ and fill $\ints{Q+1}N$. A zero offset at even depth
is graph zero, while one at odd depth is graph maximum. The latter
requires whole-graph complement. The $n=2$ residual is the explicit
one-arm center-one path, including its $m=0$ and $m=1$ leaf extensions;
the many-arm induction is not invoked at $h=n-1=1$.

In every fixed case below, the center, depth one, the tip, and each
existing center leaf use the already proved constructions of
Section~\ref{sec:centerzero}: center zero, full-graph center-one
maximum plus complement and arm permutation, the root-$R$ tip shell,
and removal/reattachment of one named leaf plus complement. Thus the
tables need to cover precisely arm depths $2,\ldots,k-1$.

\subsection{Length six: two older radius-three shells}
At $R=3$, the following words both have root $4$, each nonroot parity
inventory $\{0,1,2\}$, and adjacent-sum inventory $\ints05$:
\begin{verbatim}
H3 = 4,1,0,0,2,2,1    sums 5,1,0,2,4,3    zero indices 2,3
A3 = 4,1,2,2,0,0,1    sums 5,3,4,2,0,1    zero indices 4,5
\end{verbatim}
The selected depth equals the shell index. H3 therefore supplies depths
$2,3$, and A3 supplies $4,5$, using full complement at depths $3,5$.
H3 and A3 appeared earlier in private project records, but both are
also exact encodings of a public construction. Reversing triples of the
length-six Walecki word gives the three-twizzler
$[2,6,1,4,3,5]$ described by Hicks--Ollis--Schmitt~\cite{hos}.
Reversal and translation give $[4,2,3,0,5,1]$; complementing its
reversal gives $[4,0,5,2,3,1]$. For either word $a$, prepend root $4$
and encode odd one-based entries by $5-a_j$, even entries by $a_j$.
The results are A3 and H3, respectively.

The full fixed-$k=6$ conclusion can therefore be derived from older
public ingredients: the root-one residual obtained from Rofa's hub-zero
construction and Krop's subdivision construction~\cite{rofa,krop},
the adjacent-maximum leaf extension, the preceding public words, and
Shan--Zhong's path attachment~\cite[Theorem 2]{sz}. Ollis's
Example 5.7~\cite{ollis} also prints the one-arm path
$[1,6,0,4,3,5,2]$ used at the residual boundary.
This is an explicit deduction from older constructions, not a claim
that an earlier paper stated the entire all-$n,m$ spider theorem.
Neither the seed words nor fixed-$k=6$ coverage is claimed as an original
contribution here. The general shell transfer at its exact
lower bound $R=3$ proves all $n\ge2,m\ge0$; worldwide priority for the
full even-length family remains unknown.

\subsection{Length eight: three complete midpoint-alpha paths}
For $k=8$, each word below has length $17$, midpoint entry $8$ at
index $8$, labels $\ints0{16}$, every edge crossing cut $8$, and
weights $\ints1{16}$. The displayed weight word makes the latter
inventory fully inspectable. On the selected left arm, zero occurs at
depth $d$ and maximum $16$ at depth $d+1$:
\begin{verbatim}
d=2: 3,15,2,12,1,16,0,14,8,9,7,10,6,11,4,13,5
     12,13,10,11,15,16,14,6,1,2,3,4,5,7,9,8
d=4: 5,15,1,16,0,13,4,11,8,9,7,12,6,10,2,14,3
     10,14,15,16,13,9,7,3,1,2,5,6,4,8,12,11
d=6: 3,16,0,15,1,13,2,11,8,9,7,14,4,12,6,10,5
     13,16,15,14,12,11,9,3,1,2,7,10,8,6,4,5
\end{verbatim}
The first line of each pair is its path; the second lists consecutive
absolute edge differences. The three sources cover $2/3$, $4/5$ and
$6/7$ by direct zero and whole complement. The older radius-four
near-tip seed $[5,2,2,3,3,0,0,1,1]$ independently overlaps depths
$5/6$: both parity inventories are $\ints03$, and its adjacent sums
are $7,4,5,6,3,0,1,2$. That older seed is not needed for complete
coverage here. The $p$-core bridge requires $p\ge4$, so it is not
invoked at $k=8$.

\subsection{Length ten: one path and four inherited shells}
For $k=10$, this complete $21$-entry path has midpoint $10$ and cut
$10$, label inventory $\ints0{20}$, crossing edges and weight inventory
$\ints1{20}$:
\begin{verbatim}
7,15,5,14,2,19,1,20,0,16,10,11,9,12,8,13,6,17,4,18,3
8,10,9,12,17,18,19,20,16,6,1,2,3,4,5,7,11,13,14,15
\end{verbatim}
The first line is the path and the second its consecutive weights.
Zero is at index $8$ (depth $2$); maximum $20$ is at index $7$
(depth $3$). The inherited radius-five shell words below have root
$6$, each side inventory $\ints04$, and complete sum inventory
$\ints09$. Their listed zero indices produce the shown deficits
$c=10-j$, hence selected depths $10-c$:
\begin{verbatim}
deficits 1,2: 6,3,3,4,4,1,2,2,0,0,1  zeros 8,9
deficits 3,4: 6,3,3,4,4,1,0,0,2,2,1  zeros 6,7
deficit  5:   6,3,4,4,2,0,0,1,3,2,1  zero  5
deficits 6,7: 6,3,2,0,0,1,3,4,4,2,1  zeros 3,4
\end{verbatim}
The first is the old A3 shell lifted once by prefix
$[6,3,3,4]$; the others are old direct radius-five C, J and E
seeds. The new prefix sums, including its join, are $9,6,7,8$;
the old sums remain $\ints05$. Thus the four shells give depths
$3,\ldots,9$ and the path supplies the remaining depth $2$.
The aggregate last-eight theorem is not used at this radius.

\subsection{Length twelve: one path and four old direct shells}
For $k=12$, the complete path word and consecutive weights are
\begin{verbatim}
7,17,9,16,3,22,2,23,1,24,0,18,12,13,11,14,10,15,6,20,8,19,4,21,5
10,8,7,13,19,20,21,22,23,24,18,6,1,2,3,4,5,9,14,12,11,15,17,16
\end{verbatim}
Its labels are $\ints0{24}$, its weights $\ints1{24}$, midpoint
index $12$ carries $12$, and every edge crosses cut $12$. The zero
and maximum occur at indices $10,9$, giving depths $2,3$.
The following old radius-six shells have root $7$, both side inventories
$\ints05$, and adjacent-sum inventory $\ints0{11}$:
\begin{verbatim}
I (1,2): 7,4,3,1,2,3,5,5,4,2,0,0,1  zeros 10,11
D (3,4): 7,4,3,3,5,5,4,1,0,0,2,2,1  zeros 8,9
F (5,6): 7,4,5,5,3,1,0,0,2,3,4,2,1  zeros 6,7
G (7,8): 7,4,3,1,0,0,2,3,5,5,4,2,1  zeros 4,5
\end{verbatim}
Parentheses give deficits. Direct application at $R=6$ covers
depths $4,\ldots,11$. These are inherited near-tip seeds, not new
shell discoveries. The aggregate last-eight theorem has the stronger
sufficient domain $k\ge16$ and is not applied here.

\subsection{Length fourteen: path, two cores, and four lifted shells}
For $k=14$, a complete path source is
\begin{verbatim}
8,23,7,20,9,21,4,22,2,27,1,28,0,24,14,
15,13,16,12,17,11,18,10,19,5,26,3,25,6
\end{verbatim}
Its consecutive weights are
\begin{verbatim}
15,16,13,11,12,17,18,20,25,26,27,28,24,10,
1,2,3,4,5,6,7,8,9,14,21,23,22,19
\end{verbatim}
The labels are $\ints0{28}$, weights $\ints1{28}$, every edge
crosses cut $14$, and midpoint index $14$ is $14$. Zero at index
$12$ and maximum $28$ at index $11$ yield depths $2,3$.

At $p=5$, the two following boundary cores have root $3$, terminal
offset $p-4=1$, both parity inventories $\ints04$, and full adjacent
sum inventory $\ints08$. We print the ordered sums and zero indices:
\begin{verbatim}
C34 = 3,2,0,0,1,3,4,4,2,1  sums 5,2,0,1,4,7,8,6,3
      zero indices 2,3
C56 = 3,4,4,2,0,0,1,3,2,1  sums 7,8,6,2,0,1,4,5,3
      zero indices 4,5
\end{verbatim}
Lemma~\ref{lem:corelift} gives depths $3/4$ and $5/6$; the depth-three
overlap with the path is deliberate. The valid source contract is the
generic $p\ge4$ bridge, not its later $P(k)$ prefix theorem.

At $R=7$, the following four shells have root $8$, each nonroot side
inventory $\ints06$, full sum inventory $\ints0{13}$ and zeros at
the positions specified by their deficits:
\begin{verbatim}
1,2: 8,5,5,6,6,3,3,4,4,1,2,2,0,0,1  zeros 12,13
3,4: 8,5,5,6,6,3,3,4,4,1,0,0,2,2,1  zeros 10,11
5:   8,5,5,6,6,3,4,4,2,0,0,1,3,2,1  zero  9
6,7: 8,5,5,6,6,3,2,0,0,1,3,4,4,2,1  zeros 7,8
\end{verbatim}
They come respectively from the old A3 seed lifted twice, old
radius-five C once, old radius-five J once, and old radius-five E
once. Each lift prepends $[r+3,r,r,r+1]$ and contributes new sums
$2r,2r+1,2r+2,2r+3$ including the join, preserving the old zero
deficit. These exact $R=7$ instances cover depths $7,\ldots,13$.
No claim is made that the general eight-depth theorem works at $R=7$.

\subsection{Length sixteen: older size-six cores and radius-eight shells}
The following three pre-existing $p=6$ boundary cores have first
offset $3$, last offset $p-4=2$, both parity inventories $\ints05$,
and adjacent-sum inventory $\ints0{10}$:
\begin{verbatim}
B6_2 = 3,0,0,1,1,3,2,4,4,5,5,2
       sums 3,0,1,2,4,5,6,8,9,10,7  zero indices 1,2
B6_4 = 3,1,1,0,0,3,2,4,4,5,5,2
       sums 4,2,1,0,3,5,6,8,9,10,7  zero indices 3,4
B6_6 = 3,3,2,1,1,0,0,4,4,5,5,2
       sums 6,5,3,2,1,0,4,8,9,10,7  zero indices 5,6
\end{verbatim}
Lemma~\ref{lem:corelift} transfers these to depths $2/3$, $4/5$,
and $6/7$. The cores are verbatim from the older private
\texttt{FixedDepthSeeds} source, where their contracts were already
checked; the fixed-length composition is not a new core family.

For $R=8$, prefix $[9,6,6,7]$ to each old radius-six shell I,D,F,G:
\begin{verbatim}
I8: 9,6,6,7,7,4,3,1,2,3,5,5,4,2,0,0,1  zeros 14,15
D8: 9,6,6,7,7,4,3,3,5,5,4,1,0,0,2,2,1  zeros 12,13
F8: 9,6,6,7,7,4,5,5,3,1,0,0,2,3,4,2,1  zeros 10,11
G8: 9,6,6,7,7,4,3,1,0,0,2,3,5,5,4,2,1  zeros 8,9
\end{verbatim}
These four words have root $9$, both side inventories $\ints07$,
and full sum inventory $\ints0{15}$: the new four sums are
$15,12,13,14$ and the old sums are $\ints0{11}$. Their zero
indices give deficits $1,\ldots,8$ and hence depths $8,\ldots,15$.
The imported all-eight-depth theorem has its genuine lower bound
$R=8$ and supplies the same actual named-graph transfer. No use of
$D_8(k)$ or of the large-radius graft is involved.

\paragraph{Exhaustion and status.}
For each of the six fixed lengths, the printed certificates cover
every selected arm depth $2,\ldots,k-1$ by one of the two exact
transfers, and the endpoint constructions cover center, depth one,
tip, and every existing named unit leaf. The table records the
partition; overlapping depths cause no difficulty:
\begin{center}\begin{tabular}{c p{48mm} p{52mm}}
\toprule
$k$ & front depths & remaining internal depths\\\midrule
6 & H3: $2,3$ & A3: $4,5$\\
8 & three paths: $2,\ldots,7$ & old B overlaps $5,6$\\
10 & path: $2,3$ & radius-five shells: $3,\ldots,9$\\
12 & path: $2,3$ & radius-six shells: $4,\ldots,11$\\
14 & path: $2,3$; two cores: $3,\ldots,6$ & radius-seven shells: $7,\ldots,13$\\
16 & three cores: $2,\ldots,7$ & radius-eight shells: $8,\ldots,15$\\
\bottomrule
\end{tabular}\end{center}
All transfers retain the named midpoint at $n=2,m=0$ and the original
center leaf at $n=2,m=1$. The finite words were discovered or collected
with computer assistance, but the universal $n,m$ claims rest on the
displayed label and weight interval proofs and the earlier residual
lemmas, not on extrapolation from finite samples.

\begin{proof}[Completion of Theorem~\ref{thm:main}]
Every even $k\ge2$ is exactly $2$, $4$, or at least $6$.
Section~\ref{sec:twofour} covers the first two cases for every named
target, all $n\ge2,m\ge0$. For the third case, the preceding
$k\ge18$ proof and the six fixed lengths $6,8,10,12,14,16$ cover
the entire remaining domain. These cases exhaust the theorem.
\end{proof}

\section{Verification, reproducibility, and limitations}
\label{sec:verification}
The infinite conclusions follow from explicit bases, induction,
inventory partitions and inequalities. The proof is the exact case
split $k=2$, $k=4$, or even $k\ge6$. Section~\ref{sec:twofour}
handles the two small lengths; the six fixed lengths and the retained
$k\ge18$ proof supply the third branch. Each construction covers the
actual designated center, every named arm position and every existing
original unit leaf, allowing a different labeling for each target.

The unchanged \texttt{check.py}, \texttt{certificates.json},
\texttt{fixed\_check.py} and \texttt{fixed-certificates.json}
retain the earlier larger-length verification scope. Their already
captured results and separate mathematical QA have their own frozen
bindings. The new boundary companion reconstructs the displayed
$k=2,4$ formulas on actual named graphs, including $n=2,m=0/1$,
and tests specific mutations. Static source checks cover the exact
revision, references, citations and TeX brace/environment structure.
Finite checks corroborate transcription and do not prove unbounded
quantifiers by extrapolation.

\paragraph{Exact component formal status.}
For $k=2$, the separate copied-source replay verifies
\texttt{K2Actual.full\_actual} and its unique-zero corollary.
It freshly compiled one theorem module over 178 prior replay snapshots,
querying 1595 public theorem closures; it was not a full 179-source rebuild.
For $k=4$, the unbounded \texttt{K4Inventory.base\_gracefulness}
discharges the prior open base premise, and
\texttt{Integration4Closed.k4\_actual} covers every actual target.
Its separate replay rebuilt five copied proof modules over 178 pinned
predecessor objects, including the universal inventory and final wrapper.
For even $k\ge6$, the previously accepted replay freshly compiled its
final wrapper over 177 pinned predecessor objects and queried 1586 public
closures; it was not a fresh 178-source rebuild.

\paragraph{Unconditional all-even integration.}
The exact kernel theorem is
\texttt{Integration2Closed.all\_even\_ge\_2\_actual}:
for natural $k,n,m$ with $k\ge2$, $k\bmod2=0$, $n\ge2$ and
every $v:\texttt{SpiderVertex}\ n\ m\ k$, it gives
\[
 \exists f,\quad
 \operatorname{Graceful}(\operatorname{spiderGraph}\ n\ m\ k,
                         nk+m,f)\ \wedge\ f(v)=0.
\]
The unique-zero corollary makes every different vertex positive.
No supplied base, shell or residual proposition remains in the final
theorem. The separate integration replay freshly compiled two copied
integration sources over 184 separately replayed predecessor objects.
It is not a fresh 186-source rebuild. Its query covers 1629 public
theorem closures, each using at most \texttt{propext},
\texttt{Classical.choice} and \texttt{Quot.sound}; the public count
does not claim a separate enumeration of every generated private
declaration. It checks 57 individually named boundary targets across
$k=2,4,6$, $n=2$ and $m=0,1$, and rejects five interface mutations.
At $m=0$ it invents no leaf target. The complete universal theorem,
not those finite cases, supplies every $n,m,k$ in its domain.
The accompanying methods and source-pin files bind each component and the
final replay.\footnote{Final integration replay report SHA-256:
\texttt{e7ef4fccc268b1bd7ea3bb7a775e45fdc930e695766eaa550cca2e4076d8862f}.}

\paragraph{Complete source-package check.}
A later separate internal check copied the entire source distribution
into a new source tree and started with an empty project-object
directory. With Lean 4.34.0 and warnings treated as errors, all 186
project modules compiled freshly from those copied sources with zero
exit and no warnings. No prior project, author or incremental-replay
object was imported in that build. This full source check is distinct
from the earlier two-module incremental integration replay above.

The fresh objects passed the 1629-name public theorem audit. An
additional package-wide audit checked all 6769 theorem declarations
whose module belongs to the exact 186-module package, including 276
private declarations and generated declarations. Both inventories use
only \texttt{propext}, \texttt{Classical.choice} and \texttt{Quot.sound},
or no axioms; their counts overlap and are not to be added. The
package-wide audit also rejected any package-local axiom declaration.
The exact unrestricted theorem headers and all 57 named boundary
targets at $k=2,4,6$, $n=2$, $m=0,1$ compiled against these fresh
objects. Three substantive source mutations of edge count, parity and
arm-count premise were rejected, followed by a successful unchanged
endpoint control. Installed Lean libraries remain external toolchain
dependencies. The source invocation convention documented in the
package must be retained for its private-declaration audit.
\footnote{Separate complete source-package QA report SHA-256:
\texttt{a12575d294fec81ce6fd538d6cf27442a30a2070d0a6a33afd2a854f75b11bec}.}

This is technical verification of the exact source package and the
unconditional actual-vertex theorem. It does not establish historical
priority, external scholarly review, or PDF/layout correctness.

Formal checks rely on the pinned Lean compiler/kernel, its standard
library and the recorded imported objects. They do not establish every
prose explanation, the literature comparison, historical priority,
external scholarly review or public-release approval.
The original texts of several classical constructions remain access
gaps; accurate later restatements are identified where used. Ollis's
Lemma 5.5 is credited in that paper to Kotzig, reference [21], not to
Gv\"ozdjak. This is a correction of Ollis's attribution, not a claim
of new access to Kotzig's original proof.

AI tools assisted documented construction/proof work, coding, literal
checking, source inspection and this additive English revision. The
separate checks are internal project checks and do not constitute
external peer review. Jordi Gartner owns and publishes the work through
Beedbyte. No result for
odd $k$, $n=1$, arbitrary mixed long-arm lengths or simultaneous
multiple zeros follows from Theorem~\ref{thm:main}; no nonexistence
claim follows from failure of a restricted shell contract.
Static TeX checks are not a PDF or rendered-layout check. The accompanying
compiler record reports the actual built-in compilation outcome for
this exact source; an infrastructure failure supports no PDF claim.

\begin{thebibliography}{9}
\bibitem{patterson} Brandon J. Patterson,
\emph{New Classes of Graceful Spiders and Related Computational Results},
Master of Science thesis, Ball State University, May 2017, Section 5.4,
Conjecture 5.4.4.
\url{https://cardinalscholar.bsu.edu/server/api/core/bitstreams/b44ff232-a480-4d64-aa44-8dff4535e28c/content}.
\bibitem{rofa} Rafael I. Rofa,
\emph{0-rotatability of classes of rooted symmetric trees. Are rooted
symmetric trees 0-rotatable?}, arXiv:2312.16235v1 (2023), introduction
pp.~3--4 and summary.
\url{https://arxiv.org/pdf/2312.16235v1}.
\bibitem{hos} Jacob Hicks, M. A. Ollis and John R. Schmitt,
\emph{Distinct partial sums in cyclic groups: polynomial method and
constructive approaches}, Journal of Combinatorial Designs
\textbf{27} (2019), 369--385. DOI:10.1002/jcd.21652.
Lemma 4.5 and its proof; author manuscript:
\url{https://community.middlebury.edu/~jschmitt/papers/HicksOllisSchmitt2018.pdf}.
\bibitem{ollis} M. A. Ollis,
\emph{Sequences in dihedral groups with distinct partial products},
Australasian Journal of Combinatorics \textbf{78}(1) (2020), 35--60,
Lemma 5.5 and Theorem 5.6.
\url{https://ajc.maths.uq.edu.au/pdf/78/ajc_v78_p035.pdf}.
% V2_CONTEXT_BEGIN additional_bibliography

\bibitem{lcr} At\'{i}lio G. Luiz, C. N. Campos and R. Bruce Richter,
\emph{Some families of 0-rotatable graceful caterpillars},
Technical Report IC-17-12, Universidade Estadual de Campinas,
4 August 2017. Lemmas 3--5, Theorems 9 and 14.
\url{https://ic.unicamp.br/~reltech/2017/17-12.pdf}.
\bibitem{sz} Songling Shan and Yucheng Zhong,
\emph{Graceful Labeling of Two Families of Spiders},
arXiv:2605.14295v2 (2026), Lemmas 1--2 and Theorem 2.
\url{https://arxiv.org/html/2605.14295v2}.
\bibitem{cattell} R. Cattell,
\emph{Graceful labellings of paths}, Discrete Mathematics
\textbf{307}(24) (2007), 3161--3176.
\url{https://doi.org/10.1016/j.disc.2007.03.046}.
Original abstract/metadata and later primary restatements inspected;
complete original proof access unresolved.
\bibitem{dn} Bogdan Dumitru and Mihai Nacu,
\emph{Alternating Extremes in Graceful Labelings of Full Binary Trees
and Spider Trees}, arXiv:2607.12597v2 (2026), Section 3.2--3.3.
\url{https://arxiv.org/html/2607.12597v2}.
\bibitem{krop} Elliot Krop,
\emph{Lobsters with an almost perfect matching are graceful},
arXiv:1402.3994v1 (2014), Theorem 2.4 and the generalized
$\Delta_{+1}$ label formula.
\url{https://arxiv.org/pdf/1402.3994v1}.
% V2_CONTEXT_END additional_bibliography

\end{thebibliography}

\appendix
\section{Complete finite core certificates}\label{app:cores}
The following lists are mathematical data for the proof. Each listed size-$p$
word has its two parity inventories equal to $0,\ldots,p-1$, adjacent sums
$0,\ldots,2p-2$, and endpoints $3,p-4$. For $B_{p,d}$ the two zeros occur
at indices $d-1,d$. For $C_{p,d}$ there is the additional window
$[4,0,0,1]$ at indices $d-2,\ldots,d+1$. These conditions are the entire
finite certificate tests; they need no inference from a search outcome.
\subsection{Sizes eight through thirteen}
\[\begin{aligned}
B_{8,2}&=[3,0,0,1,1,6,5,5,7,7,6,3,2,\\
&\quad 2,4,4].
\end{aligned}\]
\[\begin{aligned}
B_{8,4}&=[3,3,5,0,0,1,1,2,2,5,4,6,6,\\
&\quad 7,7,4].
\end{aligned}\]
\[\begin{aligned}
B_{8,6}&=[3,3,5,6,1,0,0,2,2,1,4,5,7,\\
&\quad 7,6,4].
\end{aligned}\]
\[\begin{aligned}
B_{8,8}&=[3,5,2,2,4,1,1,0,0,3,7,7,6,\\
&\quad 6,5,4].
\end{aligned}\]
\[\begin{aligned}
B_{8,10}&=[3,1,2,3,5,7,7,6,1,0,0,2,4,\\
&\quad 5,6,4].
\end{aligned}\]
\[\begin{aligned}
B_{9,2}&=[3,0,0,1,1,3,2,6,8,8,7,4,6,\\
&\quad 7,5,2,4,5].
\end{aligned}\]
\[\begin{aligned}
B_{9,4}&=[3,1,1,0,0,3,2,7,8,8,6,6,7,\\
&\quad 4,4,2,5,5].
\end{aligned}\]
\[\begin{aligned}
B_{9,6}&=[3,1,2,4,1,0,0,2,7,7,8,8,5,\\
&\quad 3,4,6,6,5].
\end{aligned}\]
\[\begin{aligned}
B_{9,8}&=[3,4,8,8,7,7,6,0,0,1,1,2,2,\\
&\quad 3,5,6,4,5].
\end{aligned}\]
\[\begin{aligned}
B_{9,10}&=[3,1,2,3,5,4,6,6,1,0,0,2,4,\\
&\quad 7,7,8,8,5].
\end{aligned}\]
\[\begin{aligned}
B_{10,2}&=[3,0,0,1,1,3,2,5,9,9,8,8,7,\\
&\quad 4,5,7,6,2,4,6].
\end{aligned}\]
\[\begin{aligned}
B_{10,4}&=[3,7,5,0,0,1,1,2,2,4,4,3,6,\\
&\quad 5,9,9,8,8,7,6].
\end{aligned}\]
\[\begin{aligned}
B_{10,6}&=[3,5,1,2,2,0,0,1,4,3,8,9,9,\\
&\quad 7,7,8,5,4,6,6].
\end{aligned}\]
\[\begin{aligned}
B_{10,8}&=[3,4,7,7,5,1,1,0,0,3,2,2,6,\\
&\quad 9,9,8,8,5,4,6].
\end{aligned}\]
\[\begin{aligned}
B_{10,10}&=[3,4,2,1,4,5,7,3,1,0,0,2,6,\\
&\quad 7,8,8,9,9,5,6].
\end{aligned}\]
\[\begin{aligned}
B_{11,2}&=[3,0,0,1,1,3,2,6,10,10,9,9,8,\\
&\quad 5,7,8,6,4,5,2,4,7].
\end{aligned}\]
\[\begin{aligned}
B_{11,4}&=[3,6,5,0,0,1,1,2,2,4,4,3,7,\\
&\quad 5,10,10,9,9,8,8,6,7].
\end{aligned}\]
\[\begin{aligned}
B_{11,6}&=[3,5,1,2,2,0,0,1,4,3,7,10,10,\\
&\quad 9,9,4,5,6,6,8,8,7].
\end{aligned}\]
\[\begin{aligned}
B_{11,8}&=[3,3,1,2,6,5,2,0,0,1,4,6,10,\\
&\quad 10,9,9,8,4,5,8,7,7].
\end{aligned}\]
\[\begin{aligned}
B_{11,10}&=[3,3,1,2,7,6,5,5,2,0,0,1,4,\\
&\quad 4,8,10,10,9,6,8,9,7].
\end{aligned}\]
\[\begin{aligned}
B_{12,2}&=[3,0,0,1,1,3,2,6,11,11,10,10,9,\\
&\quad 9,7,7,8,5,6,4,5,2,4,8].
\end{aligned}\]
\[\begin{aligned}
B_{12,4}&=[3,3,5,0,0,1,1,2,2,5,4,6,6,\\
&\quad 7,9,11,11,10,8,9,10,4,7,8].
\end{aligned}\]
\[\begin{aligned}
B_{12,6}&=[3,5,1,2,2,0,0,1,4,3,8,9,10,\\
&\quad 10,11,11,7,7,9,6,6,4,5,8].
\end{aligned}\]
\[\begin{aligned}
B_{12,8}&=[3,3,1,2,6,5,2,0,0,1,4,6,7,\\
&\quad 7,10,9,9,11,11,10,5,4,8,8].
\end{aligned}\]
\[\begin{aligned}
B_{12,10}&=[3,4,8,7,6,2,1,3,2,0,0,1,5,\\
&\quad 6,4,5,9,11,11,10,7,9,10,8].
\end{aligned}\]
\[\begin{aligned}
B_{13,2}&=[3,0,0,1,1,3,2,6,12,12,11,11,10,\\
&\quad 10,9,8,8,7,7,5,6,4,5,2,4,9].
\end{aligned}\]
\[\begin{aligned}
B_{13,4}&=[3,5,5,0,0,1,1,2,2,4,7,6,6,\\
&\quad 3,4,10,12,12,11,8,10,11,9,7,8,9].
\end{aligned}\]
\[\begin{aligned}
B_{13,6}&=[3,5,1,2,2,0,0,1,4,3,8,11,12,\\
&\quad 12,10,10,11,7,9,8,7,6,6,4,5,9].
\end{aligned}\]
\[\begin{aligned}
B_{13,8}&=[3,5,8,10,9,6,4,0,0,1,1,2,5,\\
&\quad 7,10,11,11,12,12,8,6,3,2,4,7,9].
\end{aligned}\]
\[\begin{aligned}
B_{13,10}&=[3,5,8,8,10,10,9,6,4,0,0,1,1,\\
&\quad 2,5,7,7,4,2,3,6,11,11,12,12,9].
\end{aligned}\]
\subsection{Anchored sizes sixteen through twenty-one}
\[\begin{aligned}
C_{16,12}&=[3,3,2,7,7,5,8,9,6,4,4,0,0,\\
&\quad 1,1,2,5,6,10,8,12,13,14,14,15,15,\\
&\quad 11,11,13,10,9,12].
\end{aligned}\]
\[\begin{aligned}
C_{16,14}&=[3,3,2,7,10,13,15,15,14,10,8,4,4,\\
&\quad 0,0,1,1,2,5,5,6,8,11,11,9,6,\\
&\quad 7,9,12,14,13,12].
\end{aligned}\]
\[\begin{aligned}
C_{16,16}&=[3,3,2,7,10,15,15,14,14,13,13,9,7,\\
&\quad 4,4,0,0,1,1,2,5,5,8,10,11,8,\\
&\quad 6,6,9,11,12,12].
\end{aligned}\]
\[\begin{aligned}
C_{16,18}&=[3,4,2,3,9,9,7,8,11,11,14,13,8,\\
&\quad 6,5,5,4,0,0,1,1,2,6,7,10,10,\\
&\quad 13,15,15,14,12,12].
\end{aligned}\]
\[\begin{aligned}
C_{16,20}&=[3,3,2,6,9,9,12,13,14,14,15,15,11,\\
&\quad 11,13,10,10,7,4,0,0,1,1,2,5,4,\\
&\quad 6,8,8,5,7,12].
\end{aligned}\]
\[\begin{aligned}
C_{16,22}&=[3,3,2,6,6,4,7,7,8,9,9,10,14,\\
&\quad 13,13,15,15,14,11,5,4,0,0,1,1,2,\\
&\quad 5,8,12,11,10,12].
\end{aligned}\]
\[\begin{aligned}
C_{16,24}&=[3,5,8,10,13,15,15,14,12,9,7,4,2,\\
&\quad 3,6,8,11,13,14,11,9,6,4,0,0,1,\\
&\quad 1,2,5,7,10,12].
\end{aligned}\]
\[\begin{aligned}
C_{16,26}&=[3,4,2,3,5,7,8,5,6,8,9,9,12,\\
&\quad 13,14,14,15,15,11,11,13,10,10,6,4,0,\\
&\quad 0,1,1,2,7,12].
\end{aligned}\]
\[\begin{aligned}
C_{17,12}&=[3,3,2,6,9,11,13,10,8,5,4,0,0,\\
&\quad 1,1,2,5,7,10,12,15,16,16,14,14,15,\\
&\quad 11,8,6,4,7,9,12,13].
\end{aligned}\]
\[\begin{aligned}
C_{17,14}&=[3,3,2,6,9,11,14,15,13,10,8,5,4,\\
&\quad 0,0,1,1,2,5,7,10,12,15,16,16,14,\\
&\quad 12,9,7,4,6,8,11,13].
\end{aligned}\]
\[\begin{aligned}
C_{17,16}&=[3,6,8,5,7,8,11,15,16,16,14,9,9,\\
&\quad 7,4,0,0,1,1,2,5,3,2,4,6,11,\\
&\quad 10,10,12,12,13,14,15,13].
\end{aligned}\]
\[\begin{aligned}
C_{17,18}&=[3,4,2,3,5,8,11,9,9,14,14,11,10,\\
&\quad 7,8,6,4,0,0,1,1,2,7,5,6,10,\\
&\quad 12,12,15,15,16,16,13,13].
\end{aligned}\]
\[\begin{aligned}
C_{17,20}&=[3,4,2,3,7,9,10,15,16,16,14,12,12,\\
&\quad 11,11,10,8,5,4,0,0,1,1,2,6,8,\\
&\quad 9,6,5,7,13,14,15,13].
\end{aligned}\]
\[\begin{aligned}
C_{17,22}&=[3,3,2,6,8,10,11,16,16,15,15,14,14,\\
&\quad 12,10,9,6,4,7,5,4,0,0,1,1,2,\\
&\quad 5,8,9,7,13,11,12,13].
\end{aligned}\]
\[\begin{aligned}
C_{17,24}&=[3,4,2,3,7,9,12,14,16,16,15,12,10,\\
&\quad 7,8,10,13,15,14,11,9,5,4,0,0,1,\\
&\quad 1,2,6,6,5,8,11,13].
\end{aligned}\]
\[\begin{aligned}
C_{17,26}&=[3,6,9,11,14,16,16,15,13,10,8,5,5,\\
&\quad 7,10,12,15,14,12,9,7,4,2,3,4,0,\\
&\quad 0,1,1,2,6,8,11,13].
\end{aligned}\]
\[\begin{aligned}
C_{18,12}&=[3,3,2,6,9,11,13,10,8,5,4,0,0,\\
&\quad 1,1,2,5,7,10,12,14,17,17,16,16,13,\\
&\quad 15,15,12,9,7,4,6,8,11,14].
\end{aligned}\]
\[\begin{aligned}
C_{18,14}&=[3,3,2,6,9,10,11,11,12,12,8,5,4,\\
&\quad 0,0,1,1,2,5,7,10,8,6,4,7,9,\\
&\quad 16,15,14,16,17,17,15,13,13,14].
\end{aligned}\]
\[\begin{aligned}
C_{18,16}&=[3,3,2,7,10,12,14,15,15,13,12,9,7,\\
&\quad 4,4,0,0,1,1,2,5,5,8,10,13,11,\\
&\quad 9,6,6,8,11,16,16,17,17,14].
\end{aligned}\]
\[\begin{aligned}
C_{18,18}&=[3,6,9,11,13,10,8,5,7,9,12,13,15,\\
&\quad 12,10,7,4,0,0,1,1,2,5,3,2,4,\\
&\quad 6,8,11,15,14,16,16,17,17,14].
\end{aligned}\]
\[\begin{aligned}
C_{18,20}&=[3,6,8,10,13,12,10,7,9,11,15,15,16,\\
&\quad 16,17,17,12,9,4,0,0,1,1,2,5,3,\\
&\quad 2,4,6,5,7,8,11,13,14,14].
\end{aligned}\]
\[\begin{aligned}
C_{18,22}&=[3,5,9,10,8,9,15,16,16,17,17,13,14,\\
&\quad 15,13,12,11,11,10,6,4,0,0,1,1,2,\\
&\quad 5,7,6,3,2,4,7,8,12,14].
\end{aligned}\]
\[\begin{aligned}
C_{18,24}&=[3,4,2,3,7,9,12,12,10,10,13,15,16,\\
&\quad 13,14,16,17,17,15,11,8,5,4,0,0,1,\\
&\quad 1,2,6,8,9,6,5,7,11,14].
\end{aligned}\]
\[\begin{aligned}
C_{18,26}&=[3,4,2,3,7,9,12,16,15,15,17,17,16,\\
&\quad 13,14,12,10,7,8,10,13,11,9,5,4,0,\\
&\quad 0,1,1,2,6,6,5,8,11,14].
\end{aligned}\]
\[\begin{aligned}
C_{19,12}&=[3,3,2,6,9,11,13,10,8,5,4,0,0,\\
&\quad 1,1,2,5,7,10,12,14,13,12,9,7,4,\\
&\quad 6,8,11,17,18,18,16,14,15,16,17,15].
\end{aligned}\]
\[\begin{aligned}
C_{19,14}&=[3,5,8,8,6,3,2,4,7,10,9,6,4,\\
&\quad 0,0,1,1,2,5,7,11,9,15,17,18,18,\\
&\quad 16,13,12,11,10,12,14,16,17,14,13,15].
\end{aligned}\]
\[\begin{aligned}
C_{19,16}&=[3,7,11,13,18,18,17,17,16,16,14,9,8,\\
&\quad 5,4,0,0,1,1,2,5,3,2,4,7,8,\\
&\quad 6,6,10,11,9,10,12,14,15,12,13,15].
\end{aligned}\]
\[\begin{aligned}
C_{19,18}&=[3,5,8,10,11,8,6,3,2,4,7,9,14,\\
&\quad 11,9,6,4,0,0,1,1,2,5,7,10,12,\\
&\quad 12,14,13,18,18,17,17,16,16,13,15,15].
\end{aligned}\]
\[\begin{aligned}
C_{19,20}&=[3,4,2,3,5,6,9,11,14,14,16,18,18,\\
&\quad 17,12,9,7,5,4,0,0,1,1,2,8,10,\\
&\quad 13,13,11,8,6,7,10,12,15,16,17,15].
\end{aligned}\]
\[\begin{aligned}
C_{19,22}&=[3,3,2,6,9,11,13,13,16,14,14,18,18,\\
&\quad 17,17,16,15,10,8,5,4,0,0,1,1,2,\\
&\quad 5,7,10,12,11,8,6,4,7,9,12,15].
\end{aligned}\]
\[\begin{aligned}
C_{19,24}&=[3,3,2,6,7,4,6,8,8,9,14,16,18,\\
&\quad 18,17,14,12,10,9,11,10,5,4,0,0,1,\\
&\quad 1,2,5,7,11,13,16,17,15,12,13,15].
\end{aligned}\]
\[\begin{aligned}
C_{19,26}&=[3,3,2,5,8,11,13,13,17,18,18,16,16,\\
&\quad 17,14,14,15,10,7,4,5,7,9,6,4,0,\\
&\quad 0,1,1,2,6,8,10,12,11,9,12,15].
\end{aligned}\]
\[\begin{aligned}
C_{20,12}&=[3,3,2,6,9,11,13,10,8,5,4,0,0,\\
&\quad 1,1,2,5,7,10,12,14,17,18,15,12,9,\\
&\quad 7,4,6,8,11,14,15,13,17,19,19,18,16,\\
&\quad 16].
\end{aligned}\]
\[\begin{aligned}
C_{20,14}&=[3,3,2,6,9,10,7,4,6,8,8,5,4,\\
&\quad 0,0,1,1,2,5,7,11,9,13,12,14,13,\\
&\quad 19,19,18,18,17,17,16,15,15,14,10,11,12,\\
&\quad 16].
\end{aligned}\]
\[\begin{aligned}
C_{20,16}&=[3,4,2,3,8,11,14,12,10,7,9,9,5,\\
&\quad 5,4,0,0,1,1,2,6,6,7,8,12,15,\\
&\quad 17,19,19,18,13,10,11,13,15,14,16,17,18,\\
&\quad 16].
\end{aligned}\]
\[\begin{aligned}
C_{20,18}&=[3,3,2,6,9,11,14,15,18,19,19,17,13,\\
&\quad 10,8,5,4,0,0,1,1,2,5,7,10,12,\\
&\quad 15,13,11,8,6,4,7,9,12,14,17,18,16,\\
&\quad 16].
\end{aligned}\]
\[\begin{aligned}
C_{20,20}&=[3,3,2,6,9,11,14,12,10,8,11,13,15,\\
&\quad 17,12,9,8,5,4,0,0,1,1,2,5,7,\\
&\quad 7,4,6,10,13,14,16,15,19,19,18,18,17,\\
&\quad 16].
\end{aligned}\]
\[\begin{aligned}
C_{20,22}&=[3,3,2,6,9,11,14,13,15,14,19,19,18,\\
&\quad 18,17,17,13,10,8,5,4,0,0,1,1,2,\\
&\quad 5,7,10,12,12,9,7,4,6,8,11,15,16,\\
&\quad 16].
\end{aligned}\]
\[\begin{aligned}
C_{20,24}&=[3,3,2,6,10,12,15,13,11,8,9,11,14,\\
&\quad 15,18,19,19,17,13,10,8,5,4,0,0,1,\\
&\quad 1,2,5,7,7,4,6,9,12,14,17,18,16,\\
&\quad 16].
\end{aligned}\]
\[\begin{aligned}
C_{20,26}&=[3,3,2,7,10,12,15,13,11,8,6,6,9,\\
&\quad 11,14,15,17,17,16,14,12,9,7,4,4,0,\\
&\quad 0,1,1,2,5,5,8,10,13,18,18,19,19,\\
&\quad 16].
\end{aligned}\]
\[\begin{aligned}
C_{21,12}&=[3,7,10,12,15,13,11,8,6,5,4,0,0,\\
&\quad 1,1,2,5,3,2,4,8,10,13,16,20,20,\\
&\quad 19,19,18,15,17,18,16,14,12,9,7,6,9,\\
&\quad 11,14,17].
\end{aligned}\]
\[\begin{aligned}
C_{21,14}&=[3,3,2,6,9,10,7,4,6,8,8,5,4,\\
&\quad 0,0,1,1,2,5,7,11,9,15,15,16,20,\\
&\quad 20,19,19,18,17,16,18,14,14,13,13,12,10,\\
&\quad 11,12,17].
\end{aligned}\]
\[\begin{aligned}
C_{21,16}&=[3,3,2,6,9,11,14,15,18,18,13,10,8,\\
&\quad 5,4,0,0,1,1,2,5,7,10,12,15,13,\\
&\quad 11,8,6,4,7,9,12,14,16,16,19,19,20,\\
&\quad 20,17,17].
\end{aligned}\]
\[\begin{aligned}
C_{21,18}&=[3,3,2,6,9,11,10,9,13,10,7,4,6,\\
&\quad 8,8,5,4,0,0,1,1,2,5,7,11,14,\\
&\quad 17,16,12,12,14,13,16,19,20,20,18,18,19,\\
&\quad 15,15,17].
\end{aligned}\]
\[\begin{aligned}
C_{21,20}&=[3,4,2,3,5,6,9,11,11,8,6,7,10,\\
&\quad 14,12,9,7,5,4,0,0,1,1,2,8,10,\\
&\quad 13,12,16,18,18,15,14,13,17,20,20,19,19,\\
&\quad 16,15,17].
\end{aligned}\]
\[\begin{aligned}
C_{21,22}&=[3,3,2,6,9,8,10,9,11,11,14,18,19,\\
&\quad 19,20,20,16,13,8,5,4,0,0,1,1,2,\\
&\quad 5,7,7,4,6,10,13,15,15,12,12,14,17,\\
&\quad 16,18,17].
\end{aligned}\]
\[\begin{aligned}
C_{21,24}&=[3,6,8,5,7,8,9,9,10,12,15,20,20,\\
&\quad 19,19,18,18,16,17,15,13,7,4,0,0,1,\\
&\quad 1,2,5,3,2,4,6,10,11,13,16,14,12,\\
&\quad 11,14,17].
\end{aligned}\]
\[\begin{aligned}
C_{21,26}&=[3,3,2,6,9,11,13,13,14,14,20,20,19,\\
&\quad 19,18,18,17,16,16,15,15,10,8,5,4,0,\\
&\quad 0,1,1,2,5,7,10,12,11,8,6,4,7,\\
&\quad 9,12,17].
\end{aligned}\]
\subsection{Extra cores for lengths eighteen, thirty-two and thirty-four}
These are boundary cores; the second subscript is a witness identifier,
not an assertion of an anchored window. The outward zero depths are
given explicitly after each list.
\[\begin{aligned}
F_{7,2}&=[3,0,0,1,1,4,2,2,6,6,5,5,4,\\
&\quad 3].
\end{aligned}\]
Zero depths: $2,3$.
\[\begin{aligned}
F_{7,4}&=[3,1,0,0,2,4,1,2,6,6,5,5,4,\\
&\quad 3].
\end{aligned}\]
Zero depths: $3,4$.
\[\begin{aligned}
F_{7,6}&=[3,2,1,1,0,0,4,4,2,5,5,6,6,\\
&\quad 3].
\end{aligned}\]
Zero depths: $5,6$.
\[\begin{aligned}
F_{7,7}&=[3,1,2,4,1,0,0,2,6,6,5,5,4,\\
&\quad 3].
\end{aligned}\]
Zero depths: $6,7$.
\[\begin{aligned}
F_{14,12}&=[3,3,1,2,8,9,12,11,7,5,2,0,0,\\
&\quad 1,4,4,5,6,9,7,6,8,11,13,13,12,\\
&\quad 10,10].
\end{aligned}\]
Zero depths: $12,13$.
\[\begin{aligned}
F_{14,14}&=[3,3,1,2,5,4,6,11,9,5,7,6,2,\\
&\quad 0,0,1,4,7,8,8,10,9,13,13,12,12,\\
&\quad 11,10].
\end{aligned}\]
Zero depths: $14,15$.
\[\begin{aligned}
F_{15,12}&=[3,3,1,2,8,9,12,12,11,5,2,0,0,\\
&\quad 1,4,4,5,6,6,7,7,8,10,10,9,13,\\
&\quad 13,14,14,11].
\end{aligned}\]
Zero depths: $12,13$.
\[\begin{aligned}
F_{15,14}&=[3,3,1,2,5,4,6,5,7,7,9,6,2,\\
&\quad 0,0,1,4,9,8,13,12,12,14,14,13,10,\\
&\quad 10,8,11,11].
\end{aligned}\]
Zero depths: $14,15$.
\section{Complete insertion cards}\label{app:cards}
Each card below meets the side inventories and fresh edge-chain inventory
(J). Positions in $P$ begin on high parity, positions in $B,D$ on low parity.
The total new size is $q=9$ for the first six cards and $q=18$ for the last.
\subsection*{Size 9, displacement 4}
\[\begin{aligned}
P&=[3,6].
\end{aligned}\]
\[\begin{aligned}
B&=[9,4].
\end{aligned}\]
\[\begin{aligned}
D&=[1,1,2,5,3,2,4,6,5,7,7,8,8,\\
&\quad 9].
\end{aligned}\]
\subsection*{Size 9, displacement 6}
\[\begin{aligned}
P&=[3,5,6,6].
\end{aligned}\]
\[\begin{aligned}
B&=[9,4].
\end{aligned}\]
\[\begin{aligned}
D&=[1,1,2,5,4,2,3,7,7,8,8,9].
\end{aligned}\]
\subsection*{Size 9, displacement 8}
\[\begin{aligned}
P&=[3,6].
\end{aligned}\]
\[\begin{aligned}
B&=[9,7,5,8,7,4].
\end{aligned}\]
\[\begin{aligned}
D&=[1,1,2,5,3,2,4,6,8,9].
\end{aligned}\]
\subsection*{Size 9, displacement 10}
\[\begin{aligned}
P&=[3,7].
\end{aligned}\]
\[\begin{aligned}
B&=[9,9,6,5,4,2,3,4].
\end{aligned}\]
\[\begin{aligned}
D&=[1,1,2,6,8,8,5,7].
\end{aligned}\]
\subsection*{Size 9, displacement 12}
\[\begin{aligned}
P&=[3,6].
\end{aligned}\]
\[\begin{aligned}
B&=[9,8,5,6,8,7,3,2,4,4].
\end{aligned}\]
\[\begin{aligned}
D&=[1,1,2,5,7,9].
\end{aligned}\]
\subsection*{Size 9, displacement 14}
\[\begin{aligned}
P&=[3,7].
\end{aligned}\]
\[\begin{aligned}
B&=[9,9,8,8,6,7,5,6,3,2,4,4].
\end{aligned}\]
\[\begin{aligned}
D&=[1,1,2,5].
\end{aligned}\]
\subsection*{Size 18, displacement 30}
\[\begin{aligned}
P&=[3,6,11,16].
\end{aligned}\]
\[\begin{aligned}
B&=[18,18,17,17,15,16,14,15,13,13,12,12,10,\\
&\quad 9,11,10,8,8,5,5,7,7,4,2,3,4].
\end{aligned}\]
\[\begin{aligned}
D&=[1,1,2,6,9,14].
\end{aligned}\]
\section{Complete finite shell certificates}\label{app:shells}
Each radius-$q$ list has root $q+1$, both nonroot parity inventories
$0,\ldots,q-1$, and all adjacent sums $0,\ldots,2q-1$. Its zero indices
and deficits are printed to make target selection explicit.
\subsection{Near-tip parity bases}
In $J_{r,c}$ the second subscript is the originally selected deficit;
both actual zero deficits are listed. Different seed orders may cover an
overlapping deficit, which is harmless.
\[\begin{aligned}
J_{5,2}&=[6,3,2,1,3,4,4,2,0,0,1].
\end{aligned}\]
Zero indices $8,9$; deficits $2,1$.
\[\begin{aligned}
J_{5,4}&=[6,3,3,4,4,1,0,0,2,2,1].
\end{aligned}\]
Zero indices $6,7$; deficits $4,3$.
\[\begin{aligned}
J_{5,5}&=[6,3,4,4,2,0,0,1,3,2,1].
\end{aligned}\]
Zero indices $5,6$; deficits $5,4$.
\[\begin{aligned}
J_{5,6}&=[6,3,2,0,0,1,3,4,4,2,1].
\end{aligned}\]
Zero indices $3,4$; deficits $7,6$.
\[\begin{aligned}
J_{6,2}&=[7,4,3,1,2,3,5,5,4,2,0,0,1].
\end{aligned}\]
Zero indices $10,11$; deficits $2,1$.
\[\begin{aligned}
J_{6,4}&=[7,4,3,3,5,5,4,1,0,0,2,2,1].
\end{aligned}\]
Zero indices $8,9$; deficits $4,3$.
\[\begin{aligned}
J_{6,6}&=[7,4,5,5,3,1,0,0,2,3,4,2,1].
\end{aligned}\]
Zero indices $6,7$; deficits $6,5$.
\[\begin{aligned}
J_{6,8}&=[7,4,3,1,0,0,2,3,5,5,4,2,1].
\end{aligned}\]
Zero indices $4,5$; deficits $8,7$.
\[\begin{aligned}
J_{9,8}&=[10,7,6,6,8,8,7,4,3,1,0,0,2,\\
&\quad 3,5,5,4,2,1].
\end{aligned}\]
Zero indices $10,11$; deficits $8,7$.
\subsection{Finite graft sources}
\[\begin{aligned}
O_{11\mathrm a}&=[12,9,8,6,5,3,2,0,0,1,3,4,6,\\
&\quad 7,9,10,10,8,7,5,4,2,1].
\end{aligned}\]
Zero indices $7,8$; deficits $15,14$.
\[\begin{aligned}
O_{11\mathrm b}&=[12,9,10,10,8,6,5,3,2,0,0,1,3,\\
&\quad 4,6,7,9,8,7,5,4,2,1].
\end{aligned}\]
Zero indices $9,10$; deficits $13,12$.
\[\begin{aligned}
O_{12}&=[13,10,9,7,6,4,3,1,0,0,2,3,5,\\
&\quad 6,8,9,11,11,10,8,7,5,4,2,1].
\end{aligned}\]
Zero indices $8,9$; deficits $16,15$.
\[\begin{aligned}
O_{15}&=[16,13,14,14,12,10,9,7,6,4,3,1,0,\\
&\quad 0,2,3,5,6,8,9,11,12,13,11,10,8,\\
&\quad 7,5,4,2,1].
\end{aligned}\]
Zero indices $12,13$; deficits $18,17$.
\[\begin{aligned}
O_{17}&=[18,15,15,16,16,13,12,10,9,7,6,4,3,\\
&\quad 1,0,0,2,3,5,6,8,9,11,12,14,14,\\
&\quad 13,11,10,8,7,5,4,2,1].
\end{aligned}\]
Zero indices $14,15$; deficits $20,19$.
\[\begin{aligned}
O_{21}&=[22,19,20,20,18,16,17,18,19,17,15,13,12,\\
&\quad 10,9,7,6,4,3,1,0,0,2,3,5,6,\\
&\quad 8,9,11,12,14,15,16,14,13,11,10,8,7,\\
&\quad 5,4,2,1].
\end{aligned}\]
Zero indices $20,21$; deficits $22,21$.
\[\begin{aligned}
O_{9\mathrm a}&=[10,7,6,4,3,1,0,0,2,3,5,6,8,\\
&\quad 8,7,5,4,2,1].
\end{aligned}\]
Zero indices $6,7$; deficits $12,11$.
\[\begin{aligned}
O_{9\mathrm b}&=[10,7,8,8,6,4,3,1,0,0,2,3,5,\\
&\quad 6,7,5,4,2,1].
\end{aligned}\]
Zero indices $8,9$; deficits $10,9$.

\end{document}
